\subsection{Characters and Grothendieck Groups}

Denote by \(\theta_{G}\) the K-algebra homomorphism from \(K\otimes_{\mathbf{Z}}\mathrm{R}_{\mathrm{K}}(\mathrm{G})\) to \(\mathcal{Z}_{\mathrm{K}}(\mathrm{G})\) that sends \(1\otimes[\pi]\) to \(\chi_{\pi}\) for every finite-dimensional representation \(\pi\) .

PROPOSITION 8. --- a) The homomorphism \(\theta_{\mathrm{G}}\) is an isomorphism from \(\mathrm{K} \otimes_{\mathbf{Z}} \mathrm{R}_{\mathrm{K}}(\mathrm{G})\) to \(\mathcal{Z}_{\mathrm{K}}(\mathrm{G})\) .

\begin{enumerate}
\def\labelenumi{\alph{enumi})}
\setcounter{enumi}{1}
\tightlist
\item
  Suppose that K has characteristic 0. Then \(\theta_{\mathrm{G}}\) defines an isomorphism from \(\mathrm{R}_{\mathrm{K}}(\mathrm{G})\) to the subring of \(\mathcal{Z}_{\mathrm{K}}(\mathrm{G})\) consisting of the linear combinations of the characters \(\chi_{\lambda}\) for \(\lambda\) in \(\widehat{\mathrm{G}}\) with integer coefficients.
\end{enumerate}

The family \(([\lambda])_{\lambda \in \widehat{\mathbf{G}}}\) is a basis of the \(\mathbf{Z}\) -module \(\mathrm{R}_{\mathrm{K}}(\mathrm{G})\) , and the family \((\chi_{\lambda})_{\lambda \in \widehat{\mathbf{G}}}\) is a basis of the K-vector space \(\mathcal{Z}_{\mathrm{K}}(\mathrm{G})\) . Assertions a) and b) follow (VIII, p.~411, Proposition 5).

\subsection{Dimension of Simple Representations}

Denote the cardinal of the group G by n.~Let \(\pi\) be a linear representation of G in a finite-dimensional K-vector space M. For every \(g \in G\) , we have \(\pi(g)^{n} = 1_{M}\) , so the minimal polynomial of \(\pi(g)\) divides \(T^{n} - 1\) . Since \(n \cdot 1 \neq 0\) in K, this minimal polynomial is separable (V, §11, No.~2, p.~78), and since the field K is algebraically closed, the endomorphism \(\pi(g)\) of M is diagonalizable (VII, §5, No.~7, p.~40, Proposition 12). The eigenvalues of \(\pi(g)\) are n-th roots of unity, and for every \(\alpha \in K\) , the geometric multiplicity of \(\alpha\) as an eigenvalue of \(\pi(g)\) (VII, §5, No.~2, p.~30, Definition 1) is equal to the multiplicity of \(\alpha\) as a root of the characteristic polynomial of \(\pi(g)\) . Denote by \(\mathcal{O}_{n}\) the subgroup of K generated by the set \(\mu_{n}(K)\) of n-th roots of unity; it is a finitely generated Z-module and a subring of K. The character of \(\pi\) takes its values in \(\mathcal{O}_{n}\) .

PROPOSITION 9. --- Suppose that the field K has characteristic zero. Then the degree of every simple representation of G divides the cardinal of G.

Let \((\mathrm{V},\pi)\) be a simple representation of G and \(\chi\) be its character. For every element a of \(\mathrm{Z}(\mathrm{K}[\mathrm{G}])\) , the endomorphism \(\pi(a)\) of V is a homothety (VIII, p.~47, Theorem 1); denote by \(\varphi(a)\) the scalar such that \(\pi(a)=\varphi(a)_{\mathrm{V}}\) .

The resulting mapping \(\varphi\) from \(Z(K[G])\) to \(K\) is an algebra homomorphism. Let us take \(a = \sum_{g \in G} \chi(g^{-1})g\) ; by the remark of VIII, p.~408, we have \(\varphi(a) = (\dim V)^{-1}|G|\) . On the other hand, \(a\) belongs to the subring \(\mathcal{O}_n[G] \cap Z(K[G])\) of \(K[G]\) , which is a finitely generated \(\mathbf{Z}\) -module (VII, §3, p.~15, Corollary). So the element \(\varphi(a) = (\dim V)^{-1}|G|\) of \(K\) belongs to a subring of \(K\) that is a finitely generated \(\mathbf{Z}\) -module. We conclude using the following lemma.

Lemma. --- Let L be an extension of Q. Let A be a subring of L. Suppose that A is a finitely generated Z-module. We have \(A \cap Q = Z\) .

Since the Z-module \(A \cap Q\) is finitely generated, there exists a strictly positive integer N such that \(A \cap Q\) is contained in \(\frac{1}{N}Z\) . Let x be an element of \(Q - Z\) ; we write it as \(x = \frac{p}{q}\) , where p and q are mutually prime integers and \(q \geqslant 2\) . We have \(q^{N} \geqslant 2^{N} > N\) (Set Theory, III, §3, No.~6, p.~165, Theorem 2), the integers \(p^{N}\) and \(q^{N}\) are mutually prime, and consequently \(x^{N} \notin \frac{1}{N}Z\) . It follows that x does not belong to A. This concludes the proof of the lemma.

In the next subsection, we extend Proposition 9 to the case when we only assume that the characteristic of K does not divide the order of G.

\subsection{Change of Base Field}

We keep the notation of the previous subsection. Let \(K'\) be an algebraically closed field such that the element \(n \cdot 1\) of \(K'\) is not zero. The groups \(\mu_{n}(\mathrm{K})\) and \(\mu_{n}(\mathrm{K}')\) are cyclic of order n (V, §11, No.~2, p.~78, Theorem 1). Choose an isomorphism \(\varphi\) from \(\mu_{n}(\mathrm{K})\) to \(\mu_{n}(\mathrm{K}')\) . Let \(\pi\) be a linear representation of G in a finite-dimensional K-vector space, and let \(\pi'\) be a linear representation of G in a finite-dimensional \(K'\) -vector space. We say that \(\pi\) and \(\pi'\) are related (through \(\varphi\) ) if for every \(g \in G\) and \(\omega \in \mu_{n}(\mathrm{K})\) , the multiplicity of \(\omega\) as an eigenvalue of \(\pi(g)\) is equal to the multiplicity of \(\varphi(\omega)\) as an eigenvalue of \(\pi'(g)\) . When this is the case, \(\pi\) and \(\pi'\) have the same dimension, as can be seen by taking g = 1.

Let \(\pi_1\) and \(\pi_2\) (resp. \(\pi_1'\) and \(\pi_2'\) ) be linear representations of \(G\) in finite-dimensional vector spaces over \(K\) (resp. \(K'\) ). We have the following properties:

\begin{enumerate}
\def\labelenumi{\alph{enumi})}
\item
  If \(\pi_1\) is related to \(\pi_1'\) and \(\pi_2'\) , then \(\pi_1'\) and \(\pi_2'\) are isomorphic.
\item
  If \(\pi_1\) is related to \(\pi_1'\) and \(\pi_2\) to \(\pi_2'\) , then \(\pi_1 \oplus \pi_2\) is related to \(\pi_1' \oplus \pi_2'\) and \(\pi_1 \otimes \pi_2\) is related to \(\pi_1' \otimes \pi_2'\) .
\end{enumerate}

Assertion a) follows from Corollary 5 of VIII, p.~402, and assertion b) is clear.

Here, we denote by \(\mathcal{S}_{\mathrm{K}}(\mathrm{G})\) the set of classes of simple K{[}G{]}-modules, previously denoted by \(\widehat{G}\) , and we define \(\mathcal{S}_{\mathrm{K}^{\prime}}(\mathrm{G})\) likewise. The sets \(\mathcal{S}_{\mathrm{K}}(\mathrm{G})\) and \(\mathcal{S}_{\mathrm{K}^{\prime}}(\mathrm{G})\) are both finite, with cardinal the number of conjugacy classes (VIII, p.~411, Proposition 5).

PROPOSITION 10. --- There exists a unique mapping \(\varphi_{\mathrm{G}}\) from \(\mathcal{S}_{\mathrm{K}}(\mathrm{G})\) to \(\mathcal{S}_{\mathrm{K}'}(\mathrm{G})\) such that \(\lambda\) and \(\varphi_{\mathrm{G}}(\lambda)\) are related through \(\varphi\) for every \(\lambda\) in \(\mathcal{S}_{\mathrm{K}}(\mathrm{G})\) . Moreover, \(\varphi_{\mathrm{G}}\) is bijective.

The uniqueness of \(\varphi_{G}\) follows from property a) above.

\textbf{A) Suppose that the field K has characteristic 0.}

The group \(\mu_n(\mathrm{K})\) is cyclic (V, §11, No.~2, p.~78, Theorem 1); choose a generator \(\zeta\) of this group. Consider the ring homomorphism \(\rho: \mathbf{Z}[\mathrm{X}] \to \mathcal{O}_n\) that sends X to \(\zeta\) . It is surjective. The cyclotomic polynomial \(\Phi_n(\mathrm{X})\) is irreducible in \(\mathbf{Q}[\mathrm{X}]\) (V, §11, No.~5, p.~84, Theorem 2); it is therefore the minimal polynomial of \(\zeta\) over \(\mathbf{Q}\) . The polynomial \(\Phi_n\) is monic with integer coefficients (V, §11, No.~4, p.~81). Let \(\mathrm{P} \in \mathbf{Z}[\mathrm{X}]\) be a polynomial such that \(\mathrm{P}(\zeta) = 0\) ; by Euclidean division of polynomials (IV, §1, No.~6, p.~10), there exist two polynomial Q and R in \(\mathbf{Z}[\mathrm{X}]\) such that \(\mathrm{P} = \mathrm{Q}\Phi_n + \mathrm{R}\) and \(\deg(\mathrm{R}) < \deg(\Phi_n)\) . We have \(\mathrm{R}(\zeta) = 0\) , and therefore \(\mathrm{R} = 0\) because \(\Phi_n\) is the minimal polynomial of \(\zeta\) . Consequently, the kernel of \(\rho\) is the ideal \(\Phi_n\mathbf{Z}[\mathrm{X}]\) of \(\mathbf{Z}[\mathrm{X}]\) , and \(\rho\) induces a ring isomorphism from \(\mathbf{Z}[\mathrm{X}] / \Phi_n\mathbf{Z}[\mathrm{X}]\) to \(\mathcal{O}_n\) .

Set \(\zeta' = \varphi(\zeta)\) ; it is a primitive \(n\) -th root of unity in \(\mathbf{K}'\) , and we therefore have \(\Phi_n(\zeta') = 0\) (V, §11, No.~5, p.~83, Lemma 3). Consequently, there exists a homomorphism \(\varphi_0\) from the ring \(\mathcal{O}_n\) to the field \(\mathbf{K}'\) that transforms \(\zeta\) into \(\zeta'\) ; it extends the mapping \(\varphi\) from \(\mu_n(\mathrm{K})\) to \(\mu_n(\mathrm{K}')\) . Let \(\mathcal{O}\) be the subring of \(\mathrm{K}\) consisting of the elements \(\frac{a}{n^r}\) with \(a \in \mathcal{O}_n\) and \(r \in \mathbf{N}\) . Since \(n \cdot 1\) is invertible in \(\mathbf{K}'\) , the homomorphism \(\varphi_0\) extends to a homomorphism \(\varphi_1\) from \(\mathcal{O}\) to \(\mathbf{K}'\) .

We identify the algebra \(\mathcal{O}[\mathrm{G}]\) of the group \(\mathrm{G}\) over \(\mathcal{O}\) with a subring of the algebra \(\mathrm{K}[\mathrm{G}]\) and define a ring homomorphism \(\Phi\) from \(\mathcal{O}[\mathrm{G}]\) to \(\mathrm{K}'[\mathrm{G}]\) by the formula

\[
\Phi \bigg (\sum_ {g \in \mathrm{G}} a _ {g} g \bigg) = \sum_ {g \in \mathrm{G}} \varphi_ {1} (a _ {g}) g.\tag{42}
\]

Denote by \(\mathcal{C}\) the set of conjugacy classes of G. For C in \(\mathcal{C}\) , denote by \(u_{\mathrm{C}}\) the element \(\sum_{g\in \mathrm{C}}g\) of \(\mathcal{O}[\mathrm{G}]\) ; the family \((u_{\mathrm{C}})_{\mathrm{C}\in \mathcal{C}}\) is a basis over \(\mathcal{O}\) of the center \(\mathrm{Z}(\mathcal{O}[\mathrm{G}])\) of the algebra \(\mathcal{O}[\mathrm{G}]\) (VIII, p.~398). For any element \(\lambda\) of \(\mathcal{S}_{\mathrm{K}}(\mathrm{G})\) , denote by \(\chi_{\lambda}\) its character, by \(d_{\lambda}\) its dimension, and by \(e_{\lambda}\) the

element of K{[}G{]} defined by

\[
e _ {\lambda} = | \mathrm{G} | ^ {- 1} d _ {\lambda} \sum_ {g \in \mathrm{G}} \chi_ {\lambda} (g ^ {- 1}) g.\tag{43}
\]

The family \((e_{\lambda})\) is a basis over K of the center \(Z(K[G])\) of the ring K{[}G{]} (VIII, p.~408). We have

\[
e _ {\lambda} = \sum_ {\mathrm{C} \in \mathcal {C}} \alpha_ {\lambda , \mathrm{C}} u _ {\mathrm{C}}\tag{44}
\]

with

\[
\alpha_ {\lambda , \mathrm{C}} = | \mathrm{G} | ^ {- 1} d _ {\lambda} \chi_ {\lambda} (\mathrm{C} ^ {- 1}).\tag{45}
\]

For \(C \in \mathcal{C}\) and \(\lambda \in \mathcal{S}_{\mathrm{K}}(\mathrm{G})\) , set \(\beta_{\mathrm{C},\lambda} = |\mathrm{G}| d_{\lambda}^{-1} d(\mathrm{C})^{-1} \chi_{\lambda}(\mathrm{C})\) . The entries of the matrix \((\alpha_{\lambda,\mathrm{C}})\) are in \(\mathcal{O}\), and it follows from formula (31) of VIII, p.~411 that its inverse matrix is the matrix \((\beta_{\mathrm{C},\lambda})\) , whose entries are also in \(\mathcal{O}\) by Proposition 9 of VIII, p.~415. Consequently, the family \((e_{\lambda})\) is a basis of the \(\mathcal{O}\)-module \(\mathrm{Z}(\mathcal{O}[\mathrm{G}])\) .

The elements \(\Phi(u_{\mathrm{C}}) = \sum_{g\in \mathrm{C}}g\) of \(K^{\prime}[G]\) form a basis over \(K^{\prime}\) of the center \(Z(K^{\prime}[G])\) of the ring \(K^{\prime}[G]\) . We have

\[
\Phi (e _ {\lambda}) = \sum_ {\mathrm{C} \in \mathcal {C}} \varphi_ {1} (\alpha_ {\lambda , \mathrm{C}}) \Phi (u _ {\mathrm{C}}),\tag{46}
\]

and the matrix with entries \(\varphi_{1}(\alpha_{\lambda,\mathrm{C}})\) is invertible. The family of the \(\Phi(e_{\lambda})\) is therefore a basis of \(\mathrm{Z}(\mathrm{K}^{\prime}[\mathrm{G}])\) . The family \((e_{\lambda})\) is a partition of the idempotent 1 in \(\mathrm{Z}(\mathrm{K}[\mathrm{G}])\) (VIII, p.~146 and p.~408); in other words, we have

\[
\sum_ {\lambda} e _ {\lambda} = 1, \quad e _ {\lambda} ^ {2} = e _ {\lambda}, \quad e _ {\lambda} e _ {\mu} = 0 \quad \mathrm{if} \quad \lambda \neq \mu .
\]

It follows that the family of the \(\Phi(e_{\lambda})\) is a partition of the idempotent 1 in \(Z(K^{\prime}[G])\) ; since this family is a basis of \(Z(K^{\prime}[G])\) over \(K^{\prime}\) , its elements are the indecomposable idempotents in \(Z(K^{\prime}[G])\) (VIII, p.~148, Remark 4).

For \(\lambda'\) in \(\mathcal{S}_{\mathrm{K}'}(\mathrm{G})\) , define \(\chi_{\lambda'}, d_{\lambda'}\) , and \(e_{\lambda'}\) as above. By VIII, p.~408, the elements \(e_{\lambda'}\) are the indecomposable idempotents in \(\mathrm{Z}(\mathrm{K}'[\mathrm{G}])\) . Hence there exists a bijection \(\varphi_{\mathrm{G}}\) from \(\mathcal{S}_{\mathrm{K}}(\mathrm{G})\) to \(\mathcal{S}_{\mathrm{K}'}(\mathrm{G})\) such that \(\Phi(e_{\lambda}) = e_{\varphi_{\mathrm{G}}(\lambda)}\) for every \(\lambda\) in \(\mathcal{S}_{\mathrm{K}}(\mathrm{G})\) .

Let \(\lambda\) in \(\mathcal{S}_{\mathrm{K}}(\mathrm{G})\) ; set \(\lambda' = \varphi_{\mathrm{G}}(\lambda)\) . Let \((\mathrm{V}_{\lambda}, \pi_{\lambda})\) (resp. \((\mathrm{V}_{\lambda'}, \pi_{\lambda'})\) ) be a linear representation of G whose associated K{[}G{]}-module (resp. K'{[}G{]}-module) has class \(\lambda\) (resp. \(\lambda'\) ). Let us prove that \(\lambda\) and \(\lambda'\) are related. Let \(g\) be an element of G. Let \(\delta(\mathrm{T})\) be the determinant of the endomorphism \(1 + \mathrm{T}\pi_{\lambda}(g)\)

of the K{[}T{]}-module K{[}T{]} \(\otimes_{\mathrm{K}}\mathrm{V}_{\lambda}\) . Let \(\omega_{1},\ldots ,\omega_{d_{\lambda}}\) be the eigenvalues of \(\pi_{\lambda}(g)\) ; we have

\[
\delta (\mathrm{T}) = (1 + \mathrm{T} \omega_ {1}) \dots (1 + \mathrm{T} \omega_ {d _ {\lambda}}).
\]

We define \(\delta'(T)\) likewise, and we denote the eigenvalues of \(\pi_{\lambda'}(g)\) by \(\omega_1', \ldots, \omega_{d_{\lambda'}}'\) . The \(\mathcal{O}\) -module \(\mathcal{O}[G]\) is free with basis \(G\) , and the \(K\) -vector space \(K[G]\) has basis \(G\) . Denote by \(\Delta(T)\) the determinant of multiplication by \(1 + e_{\lambda}gT\) in the \(\mathcal{O}[T]\) -module \(\mathcal{O}[T] \otimes_{\mathcal{O}} \mathcal{O}[G]\) . It is also the determinant of multiplication by \(1 + e_{\lambda}gT\) in the \(K[T]\) -module \(K[T] \otimes_K K[G]\) . Let \(\overline{\varphi}_1\) be the homomorphism from \(\mathcal{O}[T]\) to \(K'[T]\) that extends \(\varphi_1\) and sends \(T\) to \(T\) . Since \(G\) is a basis of the \(K'\) -vector space \(K'[G]\) , the polynomial \(\overline{\varphi}_1(\Delta(T))\) is equal to the determinant \(\Delta'(T)\) of multiplication by \(1 + e_{\lambda'}gT\) in the \(K'\) -vector space \(K'[G]\) .

The algebra K{[}G{]} is the direct sum of its simple components \(e_{\mu}K[G]\) for \(\mu\) running through \(\mathcal{S}_{\mathrm{K}}(\mathrm{G})\) . For \(\mu\) different from \(\lambda\) , the element \(e_{\lambda}g\) annihilates \(e_{\mu}K[G]\) . Moreover, multiplication by \(e_{\lambda}g\) coincides with multiplication by g in \(e_{\lambda}K[G]\) . In view of VIII, p.~409 and Example 6 of VIII, p.~400, the representation of G in \(e_{\lambda}K[G]\) is the direct sum of \(d_{\lambda}\) representations of class \(\lambda\) . We consequently have \(\Delta(\mathrm{T}) = \delta(\mathrm{T})^{d_{\lambda}}\) .

Analogously, we have \(\Delta^{\prime}(\mathrm{T}) = \delta^{\prime}(\mathrm{T})^{d_{\lambda^{\prime}}}\) .

From the relation \(\Delta'(\mathrm{T}) = \overline{\varphi}_1(\Delta(\mathrm{T}))\) , we deduce first that \(d_{\lambda}^2 = d_{\lambda'}^2\) , and therefore \(d_{\lambda} = d_{\lambda'}\) , and then that the sequence \(\varphi(\omega_1), \ldots, \varphi(\omega_{d_\lambda})\) can be deduced from the sequence \((\omega_1', \ldots, \omega_{d_{\lambda'}'}')\) by a permutation of the set of indices.

Since this is true for every element \(g\) of \(G\) , the representations \(\lambda\) and \(\lambda'\) are related. We have therefore proved Proposition 10 when the field \(K\) has characteristic 0.

\textbf{B) General case.}

Let L be an algebraically closed field of characteristic 0 (for example, an algebraic closure of Q). Denote by \(\mathcal{S}_{\mathrm{L}}(\mathrm{G})\) the set of classes of simple L{[}G{]}-modules. Choose an isomorphism \(\eta\) from the group \(\mu_{n}(\mathrm{L})\) to the group \(\mu_{n}(\mathrm{K})\) , and set \(\eta' = \varphi \circ \eta\) . By part A) of the proof, there exist bijections

\[
\eta_ {\mathrm{G}}: \mathcal {S} _ {\mathrm{L}} (\mathrm{G}) \to \mathcal {S} _ {\mathrm{K}} (\mathrm{G}), \quad \eta_ {\mathrm{G}} ^ {\prime}: \mathcal {S} _ {\mathrm{L}} (\mathrm{G}) \to \mathcal {S} _ {\mathrm {K^ {\prime}}} (\mathrm{G})
\]

with the following property: for every \(\lambda\) in \(\mathcal{S}_{\mathrm{L}}(\mathrm{G})\) , the representations \(\lambda\) and \(\eta_{\mathrm{G}}(\lambda)\) are related through \(\eta\) , and the representations \(\lambda\) and \(\eta_{\mathrm{G}}^{\prime}(\lambda)\) are related through \(\eta^{\prime}\) . The bijection \(\varphi_{G} = \eta_{G}^{\prime} \circ \eta_{G}^{-1}\) has the desired properties.

The bijection \(\varphi_{G}\) from \(\mathcal{S}_{K}(G)\) to \(\mathcal{S}_{K'}(G)\) extends to an isomorphism, also denoted by \(\varphi_{G}\) , from the Grothendieck group \(R_{K}(G)\) to the group \(R_{K'}(G)\) .

Remark 1. --- Suppose that \(K'\) is an extension of K and that the isomorphism \(\varphi\) is the mapping \(\xi \mapsto \xi \cdot 1\) ; then the mapping \(\varphi_{G}\) is given by extension of scalars from K to \(K'\) .

COROLLARY 1. --- The mapping \(\varphi_{\mathrm{G}}\) is a ring isomorphism from \(\mathrm{R}_{\mathrm{K}}(\mathrm{G})\) to \(\mathrm{R}_{\mathrm{K}^{\prime}}(\mathrm{G})\) . For every finite-dimensional representation \(\pi\) of \(\mathrm{G}\) in a K-vector space, we have \(\varphi_{\mathrm{G}}([\pi]) = [\pi^{\prime}]\) , where \(\pi^{\prime}\) is a representation related to \(\pi\) through \(\varphi\) .

This follows from the semisimplicity of the representations of G and property b) of VIII, p.~416.

COROLLARY 2. --- The dimension of each simple representation of G divides the order of G.

This follows from Proposition 10 and Proposition 9 of VIII, p.~415.

Remarks. --- 2) Suppose that the group G is abelian. We saw in the remark of VIII, p.~414 that \(\mathcal{S}_{\mathrm{K}}(\mathrm{G})\) can be identified with the set \(\operatorname{Hom}(\mathrm{G}, \mu_n(\mathrm{K}))\) . Likewise, \(\mathcal{S}_{\mathrm{K}'}(\mathrm{G})\) can be identified with \(\operatorname{Hom}(\mathrm{G}, \mu_n(\mathrm{K}'))\) . With these identifications, the bijection \(\varphi_{\mathrm{G}}\) is simply the mapping \(\chi \mapsto \varphi \circ \chi\) .

\begin{enumerate}
\def\labelenumi{\arabic{enumi})}
\setcounter{enumi}{2}
\tightlist
\item
  Let \(\pi_1\) and \(\pi_2\) be linear representations of G in finite-dimensional vector spaces over K. For \(i = 1,2\) , let \(\pi_i'\) be a representation related to \(\pi_i\) through \(\varphi\) . We have
\end{enumerate}

\[
\dim_ {K} \operatorname{Hom} _ {K} (\pi_ {1}, \pi_ {2}) = \dim_ {K ^ {\prime}} \operatorname{Hom} _ {K ^ {\prime}} (\pi_ {1} ^ {\prime}, \pi_ {2} ^ {\prime}).\tag{47}
\]

The proof follows that of Corollary VIII, p.~410, by reducing to the case when the \(\pi_{i}\) (and therefore the \(\pi_{i}^{\prime}\) ) are simple.

\begin{enumerate}
\def\labelenumi{\arabic{enumi})}
\setcounter{enumi}{3}
\tightlist
\item
  Let H be a subgroup of G of cardinal m. The isomorphism \(\varphi\) restricts to an isomorphism from \(\mu_{m}(\mathrm{K})\) to \(\mu_{m}(\mathrm{K}^{\prime})\) and, consequently, a ring isomorphism \(\varphi_{H}\) from \(\mathrm{R}_{\mathrm{K}}(\mathrm{H})\) to \(\mathrm{R}_{\mathrm{K}^{\prime}}(\mathrm{H})\) . The following diagrams commute:
\end{enumerate}

\[
\begin{array}{c c} \mathrm {R_ {K} (G)} \xrightarrow {\mathrm {Res_ {H} ^ {G}}} \mathrm {R_ {K} (H)} & \qquad \qquad \mathrm {R_ {K} (H)} \xrightarrow {\mathrm {Ind_ {H} ^ {G}}} \mathrm {R_ {K} (G)} \\ \Biggl \downarrow \varphi_ {\mathrm{G}} & \Biggl \downarrow \varphi_ {\mathrm{H}} \\ \mathrm {R_ {K^ {\prime}} (G)} \xrightarrow {\mathrm {Res_ {H} ^ {G}}} \mathrm {R_ {K^ {\prime}} (H)}, & \qquad \qquad \mathrm {R_ {K^ {\prime}} (H)} \xrightarrow {\mathrm {Ind_ {H} ^ {G}}} \mathrm {R_ {K^ {\prime}} (G)}. \end{array}
\]

The commutativity of the first diagram is obvious, and that of the second follows from it using Frobenius reciprocity and formula (47).

\begin{enumerate}
\def\labelenumi{\arabic{enumi})}
\setcounter{enumi}{4}
\tightlist
\item
  Suppose that \(G\) is the product \(G' \times G''\) of two finite groups. We define isomorphisms \(\varphi_{G'}\) and \(\varphi_{G''}\) as in the previous example.
\end{enumerate}

We have a commutative diagram

\[
\begin{array}{c} \mathrm{R} _ {\mathrm{K}} (\mathrm{G} ^ {\prime}) \otimes_ {\mathbf {Z}} \mathrm{R} _ {\mathrm{K}} (\mathrm{G} ^ {\prime \prime}) \xrightarrow {\kappa} \mathrm{R} _ {\mathrm{K}} (\mathrm{G}) \\ \Biggl \downarrow \varphi_ {\mathrm{G} ^ {\prime}} \otimes \varphi_ {\mathrm{G} ^ {\prime \prime}} \qquad \qquad \qquad \Biggl \downarrow \varphi_ {\mathrm{G}} \\ \mathrm{R} _ {\mathrm{K} ^ {\prime}} (\mathrm{G} ^ {\prime}) \otimes_ {\mathbf {Z}} \mathrm{R} _ {\mathrm{K} ^ {\prime}} (\mathrm{G} ^ {\prime \prime}) \xrightarrow {\kappa^ {\prime}} \mathrm{R} _ {\mathrm{K} ^ {\prime}} (\mathrm{G})  , \end{array}
\]

where the isomorphisms \(\kappa\) and \(\kappa'\) are those defined in VIII, p.~406.

\subsection{Complex Linear Representations}

In this subsection, we assume that K is the field C of complex numbers.

Let \((M,\pi)\) be a linear representation of G. We say that a Hermitian form \(\Phi\) on M is invariant under G if we have

\[
\Phi (\pi (g) x, \pi (g) x ^ {\prime}) = \Phi (x, x ^ {\prime})\tag{48}
\]

for all \(x, x' \in M\) and every \(g \in G\) . This also means that for every \(g \in G\) , the automorphism \(\pi(g)\) of M is unitary with respect to \(\Phi\) .

PROPOSITION 11. --- Let \((M,\pi)\) be a finite-dimensional linear representation of G.

\begin{enumerate}
\def\labelenumi{\alph{enumi})}
\item
  There exists on M a Hermitian form that is positive, separating, and invariant under G.
\item
  Suppose that the representation \(\pi\) is simple. If \(\Phi\) and \(\Psi\) are nonzero Hermitian forms on M that are invariant under G, then there exists a real number a such that \(\Psi = a\Phi\) .
\end{enumerate}

Choose a separating positive Hermitian form on the vector space M, and denote it by \(\Phi_{0}\) . We define a separating positive Hermitian form \(\Phi\) that is invariant under G by setting

\[
\Phi (x, x ^ {\prime}) = \sum_ {g \in \mathrm{G}} \Phi_ {0} (\pi (g) x, \pi (g) x ^ {\prime})\tag{49}
\]

for \(x, x' \in M\) .

Let \(\Psi\) be a Hermitian form on M; there exists a unique endomorphism A of M such that \(\Psi(x, x') = \Phi(x, Ax')\) for \(x, x'\) in M. If, moreover, \(\Psi\) is invariant under G, then the endomorphism A commutes with the automorphism \(\pi(g)\) for \(g \in G\) . If the representation \(\pi\) is simple, then by Schur's lemma (VIII, p.~47, Theorem 1), A is a homothety and there consequently exists a complex number \(a\) such that \(\Psi = a\Phi\) . Since \(\Phi\) and \(\Psi\) are Hermitian and \(\Phi\) is nonzero, \(a\) is a real number. The proposition follows.

We endow the vector space \(\mathbf{C}[\mathrm{G}]\) of complex functions on \(\mathrm{G}\) with the Hilbert space structure whose inner product is given by

\[
\langle f | f ^ {\prime} \rangle_ {\mathrm{G}} = | \mathrm{G} | ^ {- 1} \sum_ {g \in \mathrm{G}} \overline {{f (g)}} f ^ {\prime} (g).\tag{50}
\]

For any function \(f \in C[G]\) , we denote by \(f^{*}\) the function defined by

\[
f ^ {*} (g) = \overline {{f (g ^ {- 1})}}\tag{51}
\]

for \(g \in \mathbf{G}\) ; the mapping \(f \mapsto f^*\) is a semilinear involution of \(\mathbf{C}[\mathbf{G}]\) . We also have

\[
\langle f | f ^ {\prime} \rangle_ {\mathrm{G}} = \langle f ^ {*}, f ^ {\prime} \rangle_ {\mathrm{G}}\tag{52}
\]

for \(f, f' \in \mathbf{C}[\mathrm{G}]\) , with the notation of formula (28) of VIII, p.~410. We therefore have

\[
\langle f | f ^ {\prime} \rangle_ {\mathrm{G}} = | \mathrm{G} | ^ {- 2} \tau (f ^ {*} f ^ {\prime}).\tag{53}
\]

Let \((\mathrm{M},\pi)\) be a finite-dimensional linear representation of G. We endow the vector space M with the structure of a Hilbert space for which the endomorphisms \(\pi(g)\) are unitary (Proposition 11). If we denote by \(A^{*}\) the adjoint of a endomorphism A of M for this structure, then we have \(\mathrm{Tr}(A^{*})=\overline{\mathrm{Tr}(A)}\) . For every \(g\in G\) , we have \(\pi(g^{-1})=\pi(g)^{*}\) , and therefore \(\chi_{\pi}(g^{-1})=\overline{\chi_{\pi}(g)}\) ; in other words, we have \(\chi_{\pi}=\chi_{\pi}^{*}\) . The orthogonality relation for characters (VIII, p.~410, Proposition 4) then has the form

\[
\langle \chi_ {\lambda} | \chi_ {\mu} \rangle_ {\mathrm{G}} = \delta_ {\lambda \mu}\tag{54}
\]

for \(\lambda, \mu \in \widehat{\mathbf{G}}\) . It expresses the fact that the family of the characters \((\chi_{\lambda})_{\lambda \in \widehat{\mathbf{G}}}\) of the simple representations of \(\mathbf{G}\) is an orthonormal basis of the Hilbert space \(\mathrm{Z}(\mathbf{C}[\mathrm{G}])\) of central functions.

Let \(\pi\) and \(\pi'\) be finite-dimensional linear representations of G. We have the relation \(\langle \chi_{\pi}|\chi_{\pi'}\rangle_{\mathrm{G}} = \dim_{\mathbf{C}}\mathrm{Hom}_{\mathrm{G}}(\pi, \pi')\) (VIII, p.~410, Corollary). The representation \(\pi\) is irreducible if and only if \(\langle \chi_{\pi}|\chi_{\pi}\rangle_{\mathrm{G}} = 1\) .

For every element \(\lambda\) of \(\widehat{G}\) , we endow the vector space \(V_{\lambda}\) with the structure of a Hilbert space for which the automorphisms \(\pi_{\lambda}(g)\) are unitary. We denote by \(\langle v|v'\rangle_{\lambda}\) the inner product of two elements \(v, v'\) of \(V_{\lambda}\) and by \(u^{*}\) the adjoint of an endomorphism \(u\) of \(V_{\lambda}\) . Let \(A = (A_{\lambda})_{\lambda \in \widehat{G}}\) and \(A' = (A_{\lambda}')_{\lambda \in \widehat{G}}\) be elements of \(\mathrm{F}(\widehat{\mathrm{G}})\) . We write \(\mathrm{A}^* = (\mathrm{A}_\lambda^*)_{\lambda \in \widehat{\mathrm{G}}}\) . We have \(\overline{\mathcal{F}} (a^{*}) = (\overline{\mathcal{F}} (a))^*\) for every element \(a\) of \(\mathbf{C}[\mathrm{G}]\) . Set

\[
\langle \mathrm{A} | \mathrm{A} ^ {\prime} \rangle_ {\widehat {\mathrm{G}}} = | \mathrm{G} | ^ {- 2} \widehat {\tau} (\mathrm{A} ^ {*} \mathrm{A} ^ {\prime}).\tag{55}
\]

By formula (10) of VIII, p.~407, we have

\[
\langle \mathrm{A} | \mathrm{A} ^ {\prime} \rangle_ {\widehat {\mathrm{G}}} = \frac {1}{| \mathrm{G} | ^ {2}} \sum_ {\lambda \in \widehat {\mathrm{G}}} d _ {\lambda} \operatorname{Tr} (\mathrm{A} _ {\lambda} ^ {*} \mathrm{A} _ {\lambda} ^ {\prime}).\tag{56}
\]

Since \(\hat{\tau}\circ\overline{{\mathcal{F}}}=\tau\) , formulas (53) and (55) imply that the mapping \(\overline{{\mathcal{F}}}\) is an isomorphism of Hilbert spaces from C{[}G{]} to F( \(\widehat{G}\) ).

The Schur orthogonality relations (VIII, p.~410) can be reformulated using Hilbertian inner products. Relations (21) and (24) then give the following assertions. For \(\lambda \in \widehat{\mathbf{G}}\) and \(x, x', y, y'\) in \(\mathrm{V}_{\lambda}\) , we have

\[
| \mathrm{G} | ^ {- 1} \sum_ {g \in \mathrm{G}} \overline {{\langle x | \pi_ {\lambda} (g) x ^ {\prime} \rangle}} _ {\lambda} \langle y | \pi_ {\lambda} (g) y ^ {\prime} \rangle_ {\lambda} = d _ {\lambda} ^ {- 1} \overline {{\langle x | y \rangle}} _ {\lambda} \langle x ^ {\prime} | y ^ {\prime} \rangle_ {\lambda}.\tag{57}
\]

If \(\lambda\) and \(\mu\) are two distinct elements of \(\widehat{\mathrm{G}}\) , then for \(x, x'\) in \(\mathrm{V}_{\lambda}\) and \(y, y'\) in \(\mathrm{V}_{\mu}\) , we have

\[
\sum_ {g \in \mathrm{G}} \overline {{\langle x | \pi_ {\lambda} (g) x ^ {\prime} \rangle_ {\lambda}}} \langle y | \pi_ {\mu} (g) y ^ {\prime} \rangle_ {\mu} = 0.\tag{58}
\]

For every \(\lambda \in \widehat{\mathbf{G}}\) , we choose an orthonormal basis \((e_{\lambda,i})_{1\leqslant i\leqslant d_{\lambda}}\) of \(\mathrm{V}_{\lambda}\) . For any \(g\in \mathbf{G}\) , we denote by \((\pi_{ij}^{\lambda}(g))\) the matrix of the endomorphism \(\pi_{\lambda}(g)\) of \(\mathrm{V}_{\lambda}\) with respect to this basis; we have

\[
\pi_ {i j} ^ {\lambda} (g) = \langle e _ {\lambda , i} | \pi_ {\lambda} (g) e _ {\lambda , j} \rangle_ {\lambda}.\tag{59}
\]

Since the endomorphism \(\pi_{\lambda}(g)\) is unitary, its inverse is equal to \(\pi_{\lambda}(g)^*\) , so that

\[
\overline {{\pi_ {i j} ^ {\lambda} (g)}} = \pi_ {j i} ^ {\lambda} (g ^ {- 1}).\tag{60}
\]

It then follows from formulas (22) of VIII, p.~409 and (25), p.~409 that the functions \((d_{\lambda})^{1 / 2}\pi_{ij}^{\lambda}\) , for \(\lambda \in \widehat{\mathbf{G}}\) , \(1\leqslant i\leqslant d_{\lambda}\) , \(1\leqslant j\leqslant d_{\lambda}\) , form an orthonormal basis of the Hilbert space \(\mathbf{C}[\mathrm{G}]\).*

\BourbakiExercises{Exercises}

\begin{enumerate}
\def\labelenumi{\arabic{enumi})}
\item
  Let A be a commutative ring, G a group, and H a subgroup of G of finite index. Suppose that the element \((G:H)1_{A}\) of A is invertible. Let M be an A{[}G{]}-module and N an A{[}G{]}-submodule of M. Prove that if N admits a supplementary A{[}H{]}-submodule in M, then it admits a supplementary A{[}G{]}-submodule in M (adapt the proof of Theorem 1 of VIII, p.~401). Deduce that an A{[}G{]}-module that is projective as an A{[}H{]}-module is projective.
\item
  Let A be a commutative ring and G a group.
\end{enumerate}

\begin{enumerate}
\def\labelenumi{\alph{enumi})}
\item
  Suppose that the algebra A{[}G{]} of the group G is an absolutely flat ring (VIII, p.~171, Exercise 27). Prove that the ring A is absolutely flat. Let H be a subgroup of G that admits a finite generating subset; prove that the left ideal \(I_{H}\) of A{[}G{]} (VIII, p.~22, Exercise 25) admits a supplementary in A{[}G{]}, and deduce that H is finite (observe that in the opposite case, the annihilator of \(I_{H}\) would be reduced to (0)).
\item
  Keep the assumption of a). Let g be an element of G of finite order n.~Prove that \(n1_{A}\) is invertible in A (let x be an element of A such that \((1-e_{g})x(1-e_{g})=(1-e_{g})\) ; construct an element y of A such that \(1-(1-e_{g})x=y(1+e_{g}+\cdots+e_{g^{n-1}})\) ). Deduce that the order of every finite subgroup of G is invertible in A.
\item
  Suppose that A is an absolutely flat ring and that every finite subset of G generates a finite subgroup with cardinal invertible in A. Prove that the ring A{[}G{]} is absolutely flat (reduce to the case when G is finite, and use Exercise 1).
\end{enumerate}

\begin{enumerate}
\def\labelenumi{\arabic{enumi})}
\setcounter{enumi}{2}
\tightlist
\item
  Let A be a commutative ring and G a group.
\end{enumerate}

\begin{enumerate}
\def\labelenumi{\alph{enumi})}
\item
  Suppose that the algebra A{[}G{]} of the group G is semisimple. Prove that A is semisimple, that G is finite, and that its order is invertible in A (use Exercise 2 and Exercise 25 of VIII, p.~22).
\item
  Suppose that the algebra A{[}G{]} is isomorphic to the product of a finite family of fields. Prove that every subgroup of G is normal (use a), and observe that every idempotent in A{[}G{]} is central).
\end{enumerate}

In Exercises 4 through 26 below, G is a finite group, and K is an algebraically closed field; we assume that the order of G is invertible in K.

\begin{enumerate}
\def\labelenumi{\arabic{enumi})}
\setcounter{enumi}{3}
\item
  Let \((\mathrm{M},\pi)\) be a faithful \(^{(2)}\) finite-dimensional representation of G and \((\mathrm{N},\rho)\) an irreducible representation; suppose that the character \(\chi_{\pi}\) has exactly s distinct values. Prove that there exists an integer n \textless{} s such that the K{[}G{]}-module N is isomorphic to a submodule of \(M^{\otimes n}\) (calculate \(\langle\chi_{\rho},\chi_{\pi}^{n}\rangle\) , observing that the equality \(\chi_{\pi}(g)=\dim M\) is equivalent to g=1).
\item
  Let \((\mathrm{M},\pi)\) be a representation of G.
\end{enumerate}

\begin{enumerate}
\def\labelenumi{\alph{enumi})}
\item
  Let H be a subgroup of G, S a system of representatives in G of the left cosets (mod H), and N a linear subspace of M stable under H. The representation of G in M is isomorphic to the representation induced by that of H in N if and only if we have \(M = \bigoplus_{s \in S} sN\) .
\item
  Suppose that M is the direct sum of a family \((\mathrm{M}_{i})_{i\in\mathrm{I}}\) of subspaces and that there exists a transitive action of G on I such that \(\pi(g)(\mathrm{M}_{i})\subset\mathrm{M}_{gi}\) for \(g\in G\) , \(i\in I\) . Let \(o\in I\) , and let \(G_{o}\) be its stabilizer in G. Prove that \(\pi\) is isomorphic to the representation induced by the representation of \(G_{o}\) in \(M_{o}\) . If \(\pi\) is irreducible, then the condition that the action of G on I is transitive is automatically satisfied.
\end{enumerate}

\begin{enumerate}
\def\labelenumi{\arabic{enumi})}
\setcounter{enumi}{5}
\tightlist
\item
  Let X be a finite set on which G acts. We consider the representation \(\pi\) of G in \(K^{X}\) such that \(\pi(g)e_{x}=e_{gx}\) for \(g\in G\) , \(x\in X\) (``permutation representation'' associated with X).
\end{enumerate}

\begin{enumerate}
\def\labelenumi{\alph{enumi})}
\item
  For \(g \in G\) , prove that \(\chi_{\pi}(g)\) is equal to the number of points of X fixed by g.
\item
  Prove that the multiplicity of the unit representation in \(\pi\) is equal to the number of orbits of G in X.
\item
  Assume from now on that the action of G is transitive; fix a point x of X, and denote its stabilizer by H. Prove that \(\pi\) is isomorphic to the representation induced by the unit representation of H and that \(\langle\chi_{\pi},\chi_{\pi}\rangle\) is equal to the number of orbits of H in X.
\item
  Let V be the subspace of \(K^{X}\) consisting of the vectors for which the sum of the coordinates is zero; it is stable under G, and we denote by \(\pi_{X}\) the representation of G in V deduced from \(\pi\) . Prove that \(\pi_{X}\) is irreducible if and only if the action of G is doubly transitive (I, §5, p.~143, Exercise 14).
\end{enumerate}

\begin{enumerate}
\def\labelenumi{\arabic{enumi})}
\setcounter{enumi}{6}
\item
  For any representation V of G of finite dimension, we denote by \(\varphi(\mathrm{V})\) the dimension of the subspace of V fixed by G; we extend \(\varphi\) by linearity to a Z-linear form on \(\mathrm{R}_{\mathrm{K}}(\mathrm{G})\) . Let \((\mathrm{V}_{\lambda})_{\lambda\in\widehat{\mathrm{G}}}\) be the family of irreducible representations of G, and let \(\omega\) be the element \(\sum_{\lambda\in\widehat{\mathrm{G}}}\left[\mathrm{V}_{\lambda}\right]\left[\mathrm{V}_{\lambda}^{*}\right]\) of \(\mathrm{R}_{\mathrm{K}}(\mathrm{G})\) . Prove the formula \(\varphi(\omega^{n})=\sum_{\mathrm{C}}d(\mathrm{C})^{n-1}\) for every integer \(n\geqslant0\) (apply Proposition 8 of VIII, p.~415 and the second orthogonality relation for characters).
\item
  Let H and F be subgroups of G, and let S be a subset of G such that G is the disjoint union of the double classes FsH. Let \((N,\sigma)\) be a representation of H. For \(s\in G\) , set \(H_{s}=F\cap sHs^{-1}\) , and denote by \(\sigma_{s}\) the representation of \(H_{s}\) in N defined by \(\sigma_{s}(x)=\sigma(s^{-1}xs)\) .
\end{enumerate}

\begin{enumerate}
\def\labelenumi{\alph{enumi})}
\item
  Prove that the representation \(\operatorname{Res}_{\mathrm{F}}^{\mathrm{G}}(\operatorname{Ind}_{\mathrm{H}}^{\mathrm{G}}(\sigma))\) is isomorphic to the direct sum of the representations \(\operatorname{Ind}_{\mathrm{H}_{s}}^{\mathrm{F}}(\sigma_{s})\) for \(s \in S\) . (Denote the representation \(\operatorname{Ind}_{\mathrm{H}}^{\mathrm{G}}(\sigma)\) by \((\mathrm{M}, \pi)\) . The space M can be identified with the direct sum of the \(\pi(x)\mathrm{N}\) for \(x \in G/H\) . For \(s \in S\) , let \(\mathrm{M}(s)\) be the subspace of M generated by the \(\pi(x)\mathrm{N}\) with \(x \in FsH\) . Observe that \(\mathrm{M}(s)\) is the direct sum of the subspaces \(\pi(xs)\mathrm{N}\) for \(x \in F/H_{s}\) , and deduce from Exercise 5, b) that as a K{[}F{]}-module, \(\mathrm{M}(s)\) is isomorphic to \(\operatorname{Ind}_{\mathrm{H}_{s}}^{\mathrm{F}}(\sigma_{s})\) .
\item
  Take F = H (so that we have \(H_{s} = H \cap sHs^{-1}\) ). Prove that the representation \(\operatorname{Ind}_{\mathrm{H}}^{\mathrm{G}}(\sigma)\) is irreducible if and only if \(\sigma\) is irreducible and for every \(s \in G - H\) , the supports (VIII, p.~66) of the representations \(\sigma_{s}\) and \(\operatorname{Res}_{\mathrm{H}_{s}}^{\mathrm{H}}(\sigma)\) are disjoint (``Mackey's criterion'': apply a) and Frobenius reciprocity).
\item
  Assume that H is normal; then \(\operatorname{Ind}_{\mathrm{H}}^{\mathrm{G}}(\sigma)\) is irreducible if and only if \(\sigma\) is irreducible and is not isomorphic to any of its conjugates \(\sigma \circ \operatorname{int}(g)\) for \(g \in G - H\) .
\end{enumerate}

\begin{enumerate}
\def\labelenumi{\arabic{enumi})}
\setcounter{enumi}{8}
\tightlist
\item
  Let \((\mathrm{M},\pi)\) be an irreducible representation of G and H a normal subgroup of G. a) Let \(M=\bigoplus_{i\in I}M_{i}\) be the decomposition of the K{[}H{]}-module M as a direct sum of isotypical submodules. Prove that the action of G on M permutes the submodules \(M_{i}\) transitively. Deduce that \((\mathrm{M},\pi)\) is isomorphic to the representation induced by a subgroup of G containing H (apply Exercise 5, b)).
\end{enumerate}

\begin{enumerate}
\def\labelenumi{\alph{enumi})}
\setcounter{enumi}{1}
\tightlist
\item
  Let \(\mathcal{S}\) be the support of the K{[}H{]}-module M; there exists an integer e such that M is K{[}H{]}-isomorphic to \(\sum_{\lambda\in \mathcal{S}}\lambda^{e}\) .
\end{enumerate}

\begin{enumerate}
\def\labelenumi{\arabic{enumi})}
\setcounter{enumi}{9}
\tightlist
\item
  Let \((\mathrm{M},\pi)\) be an irreducible representation of G and H a normal subgroup of G. a) Suppose that H is equal to the center of G. Prove that the degree of \(\pi\) divides the index of H in G. (For \(m \geqslant 1\) , let \(H_{m}\) be the subgroup of \(H^{m}\) consisting of the elements \((h_{1},\ldots,h_{m})\) such that \(h_{1}\ldots h_{m}=1\) . Observe that the representation \(\pi^{\times m}:G^{m}\to\operatorname{End}_{K}(M^{\otimes m})\) defines an irreducible representation of \(G^{m}/H_{m}\) , and apply Corollary 2 of Proposition 10 (VIII, p.~420) to it.)
\end{enumerate}

\begin{enumerate}
\def\labelenumi{\alph{enumi})}
\setcounter{enumi}{1}
\tightlist
\item
  Suppose that H is commutative. Prove that the degree of \(\pi\) divides (G : H). (Using Exercise 9, a), and reasoning by induction on the cardinal of G, reduce to the case when the restriction of \(\pi\) to H is isotypical. Prove that \(\pi(\mathrm{H})\) is central in \(\pi(\mathrm{G})\) , and apply a).
\end{enumerate}

\begin{enumerate}
\def\labelenumi{\arabic{enumi})}
\setcounter{enumi}{10}
\tightlist
\item
  Let A and H be subgroups of G. Suppose that A is commutative and normal and that G is the semidirect product of H and A. The group H acts on \(\widehat{\mathrm{A}} = \mathrm{Hom}(\mathrm{A}, \mathrm{K}^{*})\) ; choose a system of representatives \((\alpha_{i})_{i \in \widehat{\mathrm{A}}/\mathrm{H}}\) of the orbits of H in \(\widehat{A}\) . For \(i \in \widehat{A}/H\) , let \(H_{i}\) be the stabilizer of \(\alpha_{i}\) in H, and let \(G_{i} = AH_{i}\) . For every irreducible representation \(\sigma\) of \(H_{i}\) , define a representation \(\rho_{i,\sigma}\) of \(G_{i}\) by setting \(\rho_{i,\sigma}(ah) = \alpha_{i}(a)\sigma(h)\) for \(a \in A\) and \(h \in H\) . Prove that the representations \(\operatorname{Ind}_{G_{i}}^{G}(\rho_{i,\sigma})\) , for \(i \in \widehat{A}/H\) and \(\sigma \in \widehat{H}_{i}\) , are irreducible and pairwise nonisomorphic, and that we obtain all irreducible representations of G this way (apply Exercise 9).
\end{enumerate}

¶ 12) Let H be a subgroup of G; suppose that for every \(g \in G - H\) , we have the equality \(H \cap gHg^{-1} = \{1\}\) .

\begin{enumerate}
\def\labelenumi{\alph{enumi})}
\item
  Let u be a central function on H. Prove that the function \(\operatorname{Ind}_{\mathrm{H}}^{\mathrm{G}}(u)\) coincides with u on \(H - \{1\}\) . If v is another central function on H such that \(v(1) = 0\) , then we have \(\langle \operatorname{Ind}_{\mathrm{H}}^{\mathrm{G}}(u), \operatorname{Ind}_{\mathrm{H}}^{\mathrm{G}}(v) \rangle_{\mathrm{G}} = \langle u, v \rangle_{\mathrm{H}}\) .
\item
  Let \(\lambda \in \widehat{\mathrm{H}}\) ; we denote by \(\chi_{\lambda}^{\prime}\) the function \(\chi_{\lambda} - d_{\lambda}\) on H and by \(\theta_{\lambda}\) the function \(d_{\lambda} + \mathrm{Ind}_{\mathrm{H}}^{\mathrm{G}}(\chi_{\lambda}^{\prime})\) on G. Prove that \(\theta_{\lambda}\) is the character of an irreducible representation of G of degree \(d_{\lambda}\) whose restriction to H is isomorphic to \(\lambda\) (use a) to calculate \(\langle\theta_{\lambda},\theta_{\lambda}\rangle\) , \(\langle\theta_{\lambda},1\rangle\) , and \(\langle\mathrm{Res}_{\mathrm{H}}^{\mathrm{G}}(\theta_{\lambda}),\chi_{\lambda}\rangle\) ).
\item
  Set \(\theta = \sum_{\lambda} d_{\lambda} \theta_{\lambda}\) . Prove that we have \(\theta(h) = 0\) if \(h \in \mathrm{H} - \{1\}\) and \(\theta(g) = \theta(1) = \mathrm{Card}(\mathrm{H})\) if \(g\) does not belong to any conjugates of \(\mathrm{H}\) .
\item
  Let \(L'\) be the set of elements of G that do not belong to any conjugates of H, and set \(L = L' \cup \{1\}\) . Prove that L is a normal subgroup of G (deduce from c) that L is the kernel of a representation with character \(\theta\) .
\item
  Prove that G is the semidirect product of H and L (calculate the order of L). f) Let \(h \in H - \{1\}\) . Prove that the automorphism \(\operatorname{int}(h)\) restricted to L has no fixed point other than 1. Deduce that the order of H divides \(\operatorname{Card}(L) - 1\) .
\end{enumerate}

\(g\) ) Prove that if the subgroup H has even order, then L is commutative (apply Exercise 23, \(d\) ) of I, §6, p.~145).

\begin{enumerate}
\def\labelenumi{\arabic{enumi})}
\setcounter{enumi}{12}
\tightlist
\item
  Let \(\mathcal{H}\) be a set of subgroups of \(G\) .
\end{enumerate}

\begin{enumerate}
\def\labelenumi{\alph{enumi})}
\tightlist
\item
  Prove that the following properties are equivalent:
\end{enumerate}

\begin{enumerate}
\def\labelenumi{(\roman{enumi})}
\item
  The union of the conjugates of the subgroups belonging to \(\mathcal{H}\) is equal to G.
\item
  Every character of \(G\) is a linear combination with rational coefficients of characters induced by the characters of subgroups belonging to \(\mathcal{H}\) .
\end{enumerate}

(Deduce from Proposition 6 of VIII, p.~412 that property (ii) is equivalent to the injectivity of the restriction homomorphism \(\mathcal{Z}_{\mathrm{K}}(\mathrm{G})\to \bigoplus_{\mathrm{H}\in \mathcal{H}}\mathcal{Z}_{\mathrm{K}}(\mathrm{H}).\))

\begin{enumerate}
\def\labelenumi{\alph{enumi})}
\setcounter{enumi}{1}
\item
  Let \(\mathcal{C}\) be the set of cyclic subgroups of G; it obviously has property (i). For every cyclic subgroup C of G, we denote by \(\varphi_{\mathrm{C}}\) the function on C with value \(\operatorname{Card}(\mathrm{C})\) on the generators of C and 0 on the other elements. Prove the equality \(\operatorname{Card}(\mathrm{G}) = \sum_{\mathrm{C} \in \mathcal{C}} \operatorname{Ind}_{\mathrm{C}}^{\mathrm{G}}(\varphi_{\mathrm{C}})\) .
\item
  Let \(C \in \mathcal{C}\) . Prove that the function \(\varphi_{C}\) belongs to the subgroup \(\mathrm{R}_{\mathrm{K}}(\mathrm{G})\) of \(\mathcal{Z}_{\mathrm{K}}(\mathrm{G})\) generated by the characters (reason by induction on the order of C, using b)).
\end{enumerate}

\(d\) ) Let \(\chi\) be the character of an irreducible representation G. Prove the equality \(\chi = \mathrm{Card}(\mathrm{G})^{-1}\sum_{\mathrm{C}\in \mathcal{C}}\mathrm{Ind}_{\mathrm{C}}^{\mathrm{G}}(\varphi_{\mathrm{C}}\mathrm{Res}_{\mathrm{C}}^{\mathrm{G}}(\chi))\) , which gives an explicit formula for property (ii) of \(a\) ).

\begin{enumerate}
\def\labelenumi{\arabic{enumi})}
\setcounter{enumi}{13}
\tightlist
\item
  A representation of G is called monomial if it is the direct sum of representations induced by degree 1 representations of subgroups of G.
\end{enumerate}

\begin{enumerate}
\def\labelenumi{\alph{enumi})}
\item
  Suppose that the representations of the proper subgroups of G are monomial and that G contains a noncentral commutative normal subgroup H. Prove that every faithful irreducible representation of G is monomial (observe that the restriction to H of such a representation is not isotypical, and apply Exercise 9, a)).
\item
  Deduce that all representations of supersolvable groups (I, §6, p.~146, Exercise 26), and in particular of nilpotent groups, are monomial.
\item
  Take G to be the group \(\mathbf{SL}(2,\mathbf{F}_{3})\) . Its center Z has order 2, and the group G/Z is isomorphic to the alternating group \(\mathfrak{A}_{4}\) (II, §10, p.~422, Exercise 14, g)); in particular, G is solvable. Prove that its derived group has index 3 (cf.~I, §5, p.~142,
\end{enumerate}

Exercise 10). Using formula (8) of VIII, p.~407, deduce that G admits an irreducible representation of degree 2, and prove that this representation is not monomial.

¶ 15) Suppose that all representations of G are monomial; let us prove that G is solvable. Reasoning by induction on the order of G, we may assume that every quotient group of G different from G is solvable.

\begin{enumerate}
\def\labelenumi{\alph{enumi})}
\tightlist
\item
  If G contains two distinct minimal nontrivial normal subgroups \(H_{1}\) and \(H_{2}\) , then it is solvable (consider the canonical homomorphism from G to \(G/H_{1} \times G/H_{2}\) ).
\end{enumerate}

Assume from now on that G admits a unique minimal nontrivial normal subgroup H.

\begin{enumerate}
\def\labelenumi{\alph{enumi})}
\setcounter{enumi}{1}
\tightlist
\item
  Prove that a representation of G is faithful if its restriction to H is nontrivial. c) Let \(\pi\) be a faithful representation of minimal degree; there exist subgroups \(H_{i}\) of G and representations \(\sigma_{i}\) of \(H_{i}\) of degree 1 such that \(\pi\) is isomorphic to \(\bigoplus\operatorname{Ind}_{\operatorname{H}_{i}}^{\operatorname{G}}(\sigma_{i})\) . Let \(\pi_{0}=\bigoplus\operatorname{Ind}_{\operatorname{H}_{i}}^{\operatorname{G}}(\varepsilon_{\operatorname{H}_{i}})\) , where \(\varepsilon_{H_{i}}\) denotes the unit representation. Prove that the kernel of \(\pi_{0}\) is commutative and nontrivial (observe that \(\pi_{0}\) is irreducible); deduce that G is solvable.
\end{enumerate}

\begin{enumerate}
\def\labelenumi{\arabic{enumi})}
\setcounter{enumi}{15}
\tightlist
\item
  Suppose that the characteristic of K is zero. We say that an element x of K is integral over Z if it belongs to a subring of K that is a finitely generated Z-module or, equivalently, if the subring \(Z[x]\) of K is a finitely generated Z-module \(^{*}\) (cf.~Comm. Alg., V, §1, No.~1, p.~303).
\end{enumerate}

\begin{enumerate}
\def\labelenumi{\alph{enumi})}
\item
  Prove that the set of elements of K integral over Z is a subring of K (for x, y integral over Z, consider the subring \(Z[x, y]\) ).
\item
  An element \(x\) of \(K\) is integral over \(\mathbf{Z}\) if and only if it is a root of a monic polynomial in \(\mathbf{Z}[X]\) (when \(x\) is integral, adapt the proof of Lemma 1 of VIII, p.~81 to show that it satisfies an equation of the form \(\det(xI_m - A) = 0\) with \(A \in \mathbf{M}_m(\mathbf{Z})\) ).
\item
  If x is integral over Z, then the same holds for its conjugates \(x_{1},\ldots,x_{m}\) over Q; the element \(x_{1}\cdots x_{m}\) belongs to Z (use the lemma of VIII, p.~416).
\end{enumerate}

\begin{enumerate}
\def\labelenumi{\arabic{enumi})}
\setcounter{enumi}{16}
\tightlist
\item
  Let \(\mathrm{C}\) be a conjugacy class of \(\mathrm{G}\) and \(\pi\) an irreducible representation of \(\mathrm{G}\) of degree \(d\) .
\end{enumerate}

\begin{enumerate}
\def\labelenumi{\alph{enumi})}
\item
  Prove that the element \(d^{-1}\operatorname{Card}(\mathrm{C})\chi(\mathrm{C})\) of K is integral over Z (reason as in the proof of Proposition 9 of VIII, p.~415).
\item
  Suppose that d and \(\operatorname{Card}(\mathbf{C})\) are mutually prime integers. Prove that the element \(d^{-1}\chi(\mathbf{C})\) is integral over Z.
\end{enumerate}

* \(c\) ) Suppose, moreover, \(\mathrm{K} = \mathbf{C}\) . Prove that all conjugates of \(d^{-1}\chi (\mathrm{C})\) over \(\mathbf{Q}\) have absolute value \(\leqslant 1\) . Deduce from Exercise 16, \(c\) ) that \(|\chi (\mathrm{C})|\) is zero or equal to \(d\) . Conclude that if \(\chi (\mathrm{C})\) is not zero, then \(\pi (x)\) is a homothety for every \(x\in \mathrm{C}\) . In particular, if \(\pi\) is faithful and C is not reduced to one element, then we have \(\chi (\mathrm{C}) = 0.\ast\)

\begin{enumerate}
\def\labelenumi{\arabic{enumi})}
\setcounter{enumi}{17}
\tightlist
\item
  Suppose that the group G is simple but not cyclic.
\end{enumerate}

\begin{enumerate}
\def\labelenumi{\alph{enumi})}
\item
  Prove that there is no conjugacy class different from \(\{1\}\) whose cardinal is a power of a prime number. (Let C be a class of cardinal \(p^{r}\) with \(r \geqslant 1\) . Deduce from Exercise 17, b) that \(p^{-1}\chi_{\lambda}(1)\chi_{\lambda}(C)\) is integral over Z when \(\lambda\) is different from the unit representation, and deduce a contradiction using the second orthogonality relation for characters).
\item
  Prove that the order of G is divisible by at least three distinct prime numbers (apply a) to the conjugacy class of a central element \(x \neq 1\) in a Sylow subgroup of G).
\item
  Prove that a finite group whose order is divisible by at most two distinct prime numbers is solvable (``Burnside's theorem'': reason by induction on the cardinal of the group, using b)).
\end{enumerate}

*19) Let \(\pi\) be an irreducible complex representation of G of degree \(>1\) . Prove that there exists an element \(x\) of G such that \(\chi_{\pi}(x)=0\) . (Let \(\{\lambda_1,\ldots,\lambda_s\}\) be the orbit of \(\pi\) in \(\widehat{G}\) for the action of the group \(\operatorname{Gal}(\mathbf{C}/\mathbf{Q})\) , and let \(\chi_1,\ldots,\chi_s\) be the corresponding characters. Let \(g\in G\) ; if \(\chi(g)\neq 0\) , then prove as in Exercise 17, \(c\) ) that we have \(|\chi_1(g)\ldots\chi_s(g)|\geqslant 1\) , and therefore \(\sum_i|\chi_i(g)|^2\geqslant s\) (FRV, III, §1, No.~1, p.~93, Proposition 2). If this holds for every \(g\in G\) , then deduce from the orthogonality relation for characters that we have equality, which leads to a contradiction by taking \(g=1\) .)*

\begin{enumerate}
\def\labelenumi{\arabic{enumi})}
\setcounter{enumi}{19}
\tightlist
\item
  Let G be a finite group, H a subgroup of G, \(\sigma\) a finite-dimensional representation of H, and \(\rho\) the representation of G induced by \(\sigma\) . Let C be a conjugacy class of G; the set \(C \cap H\) is the disjoint union of a finite family \((D_{i})_{i \in I}\) of conjugacy classes of H.
\end{enumerate}

\begin{enumerate}
\def\labelenumi{\alph{enumi})}
\item
  Prove the formula \(\chi_{\rho}(\mathrm{C}) = \frac{(\mathrm{G:H})}{\mathrm{Card}(\mathrm{C})}\sum_{i\in \mathrm{I}}\mathrm{Card}(\mathrm{D}_i)\chi_\sigma (\mathrm{D}_i)\) .
\item
  If \(\sigma\) is the trivial representation, then \(\chi_{\rho}(\mathrm{C})=(\mathrm{G}:\mathrm{H})\operatorname{Card}(\mathrm{C})^{-1}\operatorname{Card}(\mathrm{C}\cap\mathrm{H})\) holds.
\end{enumerate}

¶ 21) Let \(n \in \mathbf{N}\) . Denote by \(\mathcal{D}_{n}\) the subset of \(\mathbf{N}^{n}\) consisting of the elements \((p_{1},\ldots,p_{n})\) such that \(p_{1} \geqslant \cdots \geqslant p_{n}\) ; endow it with the lexicographical order. For \(\mathbf{p} \in \mathcal{D}_{n}\) , denote by \(\mathbf{p}'\) the element of \(\mathcal{D}_{n}\) defined by \(p_{i}' = p_{i} + n - i\) . Let \(\mathbf{X} = (\mathrm{X}_{1},\ldots,\mathrm{X}_{n})\) be a sequence of variables; set \(\Delta(\mathbf{X}) = \prod_{i<j}(\mathrm{X}_{i} - \mathrm{X}_{j})\) .

\begin{enumerate}
\def\labelenumi{\alph{enumi})}
\item
  Let \(\mathrm{P} \in \mathbf{Z}[\mathbf{X}]\) be an antisymmetric polynomial, that is, a polynomial satisfying the relation \(\mathrm{P}(\mathrm{X}_{\sigma(1)}, \ldots, \mathrm{X}_{\sigma(n)}) = \varepsilon(\sigma)\mathrm{P}(\mathrm{X}_1, \ldots, \mathrm{X}_n)\) for every element \(\sigma \in \mathfrak{S}_n\) . Prove that there exists a symmetric polynomial \(\mathrm{Q} \in \mathbf{Z}[\mathbf{X}]\) such that \(\mathrm{P} = \Delta\mathrm{Q}\) .
\item
  For every \(\mathbf{p} \in \mathcal{D}_n\) , set \(S_{\mathbf{p}}(\mathbf{X}) = \Delta(\mathbf{X})^{-1} \det(X_j^{p_i'})_{1 \leqslant i,j \leqslant n}\) . Prove that \(S_{\mathbf{p}}\) is a homogeneous symmetric polynomial of degree \(\sum p_i\) with integer coefficients.
\item
  If \(\mathbf{q} = (q_i)_{1\leqslant i\leqslant n}\) is an element of \(\mathcal{D}_n\) and \(P\) is an element of \(\mathbf{Q}[\mathbf{X}]\) , then we denote by \([P]_{\mathbf{q}}\) the coefficient of \(\mathbf{x}^{\mathbf{q}}\) in \(P\) . Set \(K_{\mathbf{pq}} = [S_{\mathbf{p}}]_{\mathbf{q}}\) . Prove that we have \(K_{\mathbf{pp}} = 1\) and \(K_{\mathbf{pq}} = 0\) for \(\mathbf{q} > \mathbf{p}\) (expand the equality \(\det(X_j^{p_i'}) = S_{\mathbf{p}}(\mathbf{X})\Delta(\mathbf{X})\) in the lexicographical order). Deduce that the polynomials \(S_{p}\) for \(p \in \mathcal{D}_{n}\) form a basis of the Z-submodule of Z{[}X{]} consisting of the symmetric polynomials.
\item
  Let \(\mathbf{Y} = (\mathrm{Y}_1, \ldots, \mathrm{Y}_n)\) be a sequence of variables. Prove the equality of formal power series
\end{enumerate}

\[
\prod_ {1 \leqslant i, j \leqslant n} (1 - \mathrm{X} _ {i} \mathrm{Y} _ {j}) ^ {- 1} = \sum_ {\mathbf {p} \in \mathcal {D} _ {n}} \mathrm{S} _ {\mathbf {p}} (\mathbf {X})   \mathrm{S} _ {\mathbf {p}} (\mathbf {Y})
\]

(use the equality \(\det\big((1 - X_iY_j)^{-1}\big)_{1\leqslant i,j\leqslant n} = \Delta(\mathbf{X})\Delta(\mathbf{Y})\prod_{1\leqslant i,j\leqslant n}(1 - X_iY_j)^{-1}\) (cf.~III, §8, p.~639, Exercise 5) by calculating the left-hand side using the expansion of \((1 - T)^{-1}\) as a formal power series).

\begin{enumerate}
\def\labelenumi{\alph{enumi})}
\setcounter{enumi}{4}
\tightlist
\item
  Let \(\mathrm{P} \in \mathbf{Q}[\mathbf{X}]\) be a symmetric polynomial. For every \(\mathbf{p} \in \mathcal{D}_n\) , set \(\omega_{\mathbf{p}}(\mathrm{P}) = [\mathrm{P}\Delta]_{\mathbf{p}'}\) . Prove the equality \(\omega_{\mathbf{p}}(\mathrm{S}_{\mathbf{q}}) = \delta_{\mathbf{pq}}\) . Deduce that for every \(\mathbf{q} \in \mathcal{D}_n\) , we have \([\mathrm{P}]_{\mathbf{q}} = \sum_{\mathbf{p} \in \mathcal{D}_n} \mathrm{K}_{\mathbf{pq}} \omega_{\mathbf{p}}(\mathrm{P})\) .
\end{enumerate}

\(f)\) For \(\boldsymbol{\alpha}\in\mathbf{N}^{n}\) , set \(f^{\boldsymbol{\alpha}}(\mathbf{X})=\prod_{j=1}^{n}(\sum_{i}\mathrm{X}_{i}^{j})^{\alpha_{j}}\) and \(c_{\boldsymbol{\alpha}}=\prod_{j=1}^{n}\frac{1}{j^{\alpha_{j}}\alpha_{j}!}\) . Prove that we have \(\sum_{\boldsymbol{\alpha}\in\mathbf{N}^{n}}c_{\boldsymbol{\alpha}}\omega_{\mathbf{p}}(f^{\boldsymbol{\alpha}})\omega_{\mathbf{q}}(f^{\boldsymbol{\alpha}})=\delta_{\mathbf{pq}}\) for \(\mathbf{p},\mathbf{q}\) in \(\mathcal{D}_{n}\) (observe that we have

\[
\prod_ {1 \leqslant i, j \leqslant n} (1 - \mathrm{X} _ {i} \mathrm{Y} _ {j}) ^ {- 1} = \prod_ {k = 0} ^ {\infty} \exp \bigl (\frac {1}{k} (\sum_ {i} \mathrm{X} _ {i} ^ {k}) (\sum_ {j} \mathrm{Y} _ {j} ^ {k}) \bigr) = \sum_ {\boldsymbol {\alpha} \in \mathbf {N} ^ {n}} c _ {\boldsymbol {\alpha}} f ^ {\boldsymbol {\alpha}} (\mathbf {X}) f ^ {\boldsymbol {\alpha}} (\mathbf {Y}) \bigr).
\]

¶ 22) Let n be a natural number. Denote by \(\mathcal{S}_{n}\) the set of decreasing finite sequences of strictly positive integers with sum n.~Let \(\mathbf{p} = (p_{i})_{1 \leqslant i \leqslant d}\) be an element of \(\mathcal{S}_{n}\) . For \(1 \leqslant i \leqslant d\) , denote by \(P_{i}\) the set of integers r such that \(p_{1} + \cdots + p_{i-1} < r \leqslant p_{1} + \cdots + p_{i}\) , and for \(1 \leqslant j \leqslant p_{1}\) , denote by \(P_{j}'\) the set of integers of the form \(p_{1} + \cdots + p_{i-1} + j\) with \(1 \leqslant i \leqslant d\) and \(p_{i} \geqslant j\) . The subsets \(P_{i}\) , on the one hand, and the subsets \(P_{j}'\) , on the other, form a partition of the interval {[}1, n{]} in N.

(We often represent this situation by a diagram, called a Young diagram: in the example below, we have n = 9 and \(\mathbf{p} = (4, 2, 2, 1)\) . The set \(P_{i}\) consists of the elements of the i-th row, and \(P_{j}^{\prime}\) of those of the j-th column.)

1

2

3

4

5

6

7

8

9

Let \(\mathcal{P}\) (resp. \(\mathcal{P}'\) ) be the subgroup of \(\mathfrak{S}_n\) consisting of the permutations \(\sigma\) such that \(\sigma(\mathrm{P}_i) \subset \mathrm{P}_i\) for every \(i\) (resp. \(\sigma(\mathrm{P}_j') \subset \mathrm{P}_j'\) for every \(j\) ).

\begin{enumerate}
\def\labelenumi{\alph{enumi})}
\item
  Prove that we have \(\mathcal{P} \cap \mathcal{P}' = \{1\}\) .
\item
  Let \(\sigma \in \mathfrak{S}_n\) . Show that if \(\mathrm{Card}(\mathrm{P}_i \cap \sigma(\mathrm{P}_j')) \leqslant 1\) for every \(i\) and \(j\) , then \(\sigma\) belongs to \(\mathcal{PP}'\) . Deduce that if \(\sigma \notin \mathcal{PP}'\) , then the subgroup \(\mathcal{P} \cap \sigma \mathcal{P}'\sigma^{-1}\) contains a transposition.
\item
  In the group algebra \(\mathbf{Z}[\mathfrak{S}_n]\) , set
\end{enumerate}

\[
a _ {\mathbf {p}} = \sum_ {\sigma \in \mathcal {P}} e _ {\sigma}, b _ {\mathbf {p}} = \sum_ {\tau \in \mathcal {P} ^ {\prime}} \varepsilon (\tau) e _ {\tau}, c _ {\mathbf {p}} = a _ {\mathbf {p}} b _ {\mathbf {p}}.
\]

Prove that \(c_{p}\) satisfies \(e_{\sigma}c_{p}e_{\sigma'}=\varepsilon(\sigma')c_{p}\) for every \(\sigma\in \mathcal{P}\) and \(\sigma'\in \mathcal{P}'\) and that every element that has this property belongs to \(Zc_{p}\) (write such a element as \(\sum n_{\sigma}e_{\sigma}\) , and deduce from b) that \(n_{\sigma}\) is zero for \(\sigma\notin \mathcal{P}\mathcal{P}'\) ). Deduce that for every element x of \(Z[\mathfrak{S}_n]\) , the element \(c_{p}xc_{p}\) belongs to \(Zc_{p}\) .

\(d\) ) Denote by A the \(\mathbf{Q}\) -algebra \(\mathbf{Q}[\mathfrak{S}_n]\) . Prove that the ideal \(V_{\mathbf{p}} = Ac_{\mathbf{p}}\) of A is an absolutely simple A-module. (Observe that we have \(c_{\mathbf{p}}V_{\mathbf{p}} \subset \mathbf{Q}c_{\mathbf{p}}\) by \(c\) ). Prove that if \(\mathfrak{a}\) is a left ideal strictly contained in \(V_{\mathbf{p}}\) , then we have \(c_{\mathbf{p}}\mathfrak{a} = 0\) , and therefore \(\mathfrak{a}^2 = 0\) and finally \(\mathfrak{a} = 0\) .)

\begin{enumerate}
\def\labelenumi{\alph{enumi})}
\setcounter{enumi}{4}
\tightlist
\item
  Prove that we have \(c_{\mathbf{p}}^{2} = (n!/[V_{\mathbf{p}} : \mathbf{Q}]) c_{\mathbf{p}}\) (calculate the trace of the endomorphism \(x \mapsto xc_{p}\) of A).
\end{enumerate}

¶ 23) Keep the notation of the previous exercise. Let \(\mathbf{p} = (p_i)\) and \(\mathbf{q} = (q_j)\) be two elements of \(S_n\) .

\begin{enumerate}
\def\labelenumi{\alph{enumi})}
\item
  Suppose that we have \(\mathbf{p} > \mathbf{q}\) for the lexicographical order. Prove that we have \(a_{\mathbf{p}}xb_{\mathbf{q}} = 0\) for every \(x \in \mathrm{A}\) . (Reduce to the case \(x = 1\) . Observe that if \((\mathrm{P}_i), (\mathrm{P}_j')\) denote the partitions associated with \(\mathbf{p}\) and \((\mathrm{Q}_k), (\mathrm{Q}_l')\) those associated with \(\mathbf{q}\) , then, denoting by \(s\) the smallest integer \(n\) such that \(p_n \neq q_n\) , we have \(\operatorname{Card}(\mathrm{P}_s \cap \mathrm{Q}_1') \geqslant 2\) , and construct a transposition \(\tau\) such that \(a_{\mathbf{p}}\tau = a_{\mathbf{p}}\) and \(\tau b_{\mathbf{q}} = -b_{\mathbf{q}}\) .)
\item
  Deduce that if \(p \neq q\) , then the A-modules \(V_{p}\) and \(V_{q}\) are not isomorphic (apply a), using the antiautomorphism of A that sends \(e_{\sigma}\) to \(e_{\sigma^{-1}}\) .
\item
  Prove that the mapping that sends a sequence \(\mathbf{p}\) to the representation \((\mathrm{V}_{\mathbf{p}})_{(\mathrm{K})}\) defines a bijection from \(\mathcal{S}_n\) to the set \(\mathcal{S}_{\mathrm{K}}(\mathfrak{S}_n)\) of classes of simple \(\mathrm{K}[\mathfrak{S}_n]\) -modules. d) For \(\mathbf{p} \in \mathcal{S}_n\) , denote by \(\overline{\mathbf{p}}\) the sequence defined by \(\overline{p}_j = \mathrm{Card}(\mathrm{P}_j')\) (Exercise 22). Prove that the mapping \(\mathbf{p} \mapsto \overline{\mathbf{p}}\) is an involution of \(\mathcal{S}_n\) ; the A-module \(\mathrm{V}_{\overline{\mathbf{p}}}\) is isomorphic to \(\mathrm{V}_{\mathbf{p}} \otimes_{\mathbf{Q}} \mathbf{Q}_{\varepsilon}\) , where \(\mathbf{Q}_{\varepsilon}\) denotes the 1-dimensional representation over \(\mathbf{Q}\) associated with the signature.
\item
  Describe the representations corresponding to the sequences \((1, \ldots, 1)\) and \((n)\) , as well as to the sequence \((n - 1, 1)\) .
\end{enumerate}

¶ 24) Keep the notation of Exercises 21 through 23. Let \(\mathbf{p} = (p_i)_{1 \leqslant i \leqslant d} \in \mathcal{S}_n\) ; denote by \(\mathbf{p}\) the element \((p_1, \ldots, p_d, 0, \ldots, 0)\) of \(\mathcal{D}_n\) . If \(\boldsymbol{\alpha} = (\alpha_1, \ldots, \alpha_n)\) is an element of \(\mathbf{N}^n\) such that \(\sum_i i\alpha_i = n\) , then we denote by \(C_{\boldsymbol{\alpha}}\) the conjugacy class of \(\mathfrak{S}_n\) consisting of the elements that are the product of a family of cycles with disjoint support consisting of \(\alpha_2\) cycles of order 2, \(\alpha_3\) cycles of order 3, etc.

\begin{enumerate}
\def\labelenumi{\alph{enumi})}
\item
  Prove that the representation \(\rho_{\mathbf{p}}\) of \(\mathfrak{S}_n\) in the A-module \(\mathrm{Aa}_{\mathbf{p}}\) is the representation induced by the trivial representation of \(\mathcal{P}\) . Use Exercise 20 to deduce the formula \(\chi_{\rho_{\mathbf{p}}}(\mathrm{C}_{\alpha}) = [f^{\alpha}]_{\mathbf{p}}\) .
\item
  For any \(\mathbf{q} \in \mathcal{S}_n\) , denote by \(\lambda_{\mathbf{q}}\) the representation of \(\mathfrak{S}_n\) in the A-module \(V_{\mathbf{q}}\) . Prove that there exist positive integers \(n_{\mathbf{pq}}\) such that we have \(\chi_{\rho_{\mathbf{p}}} = \sum_{\mathbf{q} \in \mathcal{S}_n} n_{\mathbf{pq}} \chi_{\lambda_{\mathbf{q}}}\) with \(n_{\mathbf{pp}} > 0\) .
\item
  Let \(\eta_{\mathbf{p}}\) be the central function on \(\mathfrak{S}_n\) with value \(\omega_{\mathbf{p}}(f^{\alpha})\) on the elements of the class \(\mathrm{C}_{\alpha}\) . Prove that \(\eta_{\mathbf{p}}\) is a linear combination of the characters \(\chi_{\lambda_{\mathbf{q}}}\) for \(\mathbf{q} \in \mathcal{S}_n\) with integer coefficients (use Exercise 21, c) and e)). Deduce that \(\eta_{\mathbf{p}}\) is equal, up to a sign, to one of the characters \(\chi_{\lambda_{\mathbf{q}}}\) . (Calculate \(\langle \eta_{\mathbf{p}}, \eta_{\mathbf{p}} \rangle\) using Exercise 21, f). d) Deduce the equality \(\chi_{\lambda_{\mathbf{p}}}(\mathrm{C}_{\alpha}) = [\Delta f^{\alpha}]_{\mathbf{p}'}\) (``Frobenius formula'').
\item
  Prove that the multiplicity of the irreducible representation \(\lambda_{p}\) in the representation \(\rho_{q}\) is equal to \(K_{pq}\) .
\end{enumerate}

\(f\) ) Prove that the degree of \(\lambda_{\mathbf{p}}\) is \(\frac{n!}{\mathbf{p}^{\prime}!}\Delta (\mathbf{p}^{\prime})\) (expand the polynomial \(\Delta (\mathbf{X})(\sum \mathrm{X}_i)^n\) ).

¶ 25) Consider the representation \(\pi_{\mathbf{n}}\) of \(\mathfrak{S}_n\) associated with the action of \(\mathfrak{S}_n\) on the set \(\mathbf{n} = [1, n]\) (Exercise 6). Let \(k\) be a positive integer and \(\chi_k\) the character of the representation \(\wedge^k \pi_{\mathbf{n}}\) .

\begin{enumerate}
\def\labelenumi{\alph{enumi})}
\item
  Let \(\sigma \in \mathfrak{S}_n\) . Prove the formula \(\chi_k(\sigma) = \sum \varepsilon (\sigma |S)\) , where the sum is taken over the set of subsets \(S\) of \(\mathbf{n}\) with \(k\) elements such that \(\sigma (S) = S\) .
\item
  Prove that \(\wedge^{k}\pi_{n}\) is the irreducible representation of \(\mathfrak{S}_n\) associated with the element \((n-k,1,\ldots,1)\) of \(\mathcal{S}_n\) (calculate the character of this representation using the Frobenius formula, observing that the only terms of the sum
\end{enumerate}

\[
\Delta (\mathbf {X}) f ^ {\boldsymbol {\alpha}} (\mathbf {X}) = \sum_ {\sigma \in \mathfrak {S} _ {n}} \varepsilon (\sigma) X _ {1} ^ {\sigma (n) - 1} \dots X _ {n} ^ {\sigma (1) - 1} f ^ {\boldsymbol {\alpha}} (\mathbf {X})
\]

for which the coefficient of \(X^{p'}\) is nonzero correspond to the permutations \(\sigma\) that satisfy \(\sigma(i) \leqslant i + 1\) for every i and \(\sigma(i) = i\) for \(1 \leqslant i < n - k\) .

¶ 26) Let q be a power of a prime number and F a finite field with q elements; take G to be the group \(\mathbf{GL}(2, \mathbf{F})\) .

\begin{enumerate}
\def\labelenumi{\alph{enumi})}
\item
  Describe the conjugacy classes of G (distinguish four types of classes, containing, respectively, 1, \(q^{2}-1\) , \(q^{2}+q\) , and \(q^{2}-q\) elements).
\item
  The group G acts on the projective line P over F; consider the representation \(\pi_{P}\) of G (Exercise 6). For any group homomorphism \(\lambda : F^{*} \to K^{*}\) , denote by \(\delta_{\lambda}\) the representation of degree 1 associated with the homomorphism \(\lambda \circ det\) from G to \(K^{*}\) , and set \(\pi_{\lambda} = \pi_{P} \otimes \delta_{\lambda}\) . Prove that \(\pi_{\lambda}\) is irreducible, and calculate its character.
\item
  Let B be the subgroup of G consisting of the upper triangular matrices. For \(\lambda,\mu\) in \(\mathrm{Hom}(\mathrm{F}^{*},\mathrm{K}^{*})\) , denote by \(\beta_{\lambda,\mu}\) the representation of B of degree 1 defined by \(\beta_{\lambda,\mu}((a_{ij}))=\lambda(a_{11})\mu(a_{22})\) , and set \(\pi_{\lambda,\mu}=\mathrm{Ind}_{\mathrm{B}}^{\mathrm{G}}(\beta_{\lambda,\mu})\) . Prove that \(\pi_{\lambda,\mu}\) is irreducible for \(\lambda\neq\mu\) , while \(\pi_{\lambda,\lambda}\) is isomorphic to \(\pi_{P}\oplus\delta_{\lambda}\) (apply Exercise 8). The representations \(\pi_{\lambda,\mu}\) and \(\pi_{\lambda',\mu'}\) are isomorphic if and only if the subsets \(\{\lambda,\mu\}\) and \(\{\lambda',\mu'\}\) of \(\mathrm{Hom}(\mathrm{F}^{*},\mathrm{K}^{*})\) are equal.
\end{enumerate}

\(d\) ) Let \(\mathrm{F}^{\prime}\) be an extension of F of degree 2; choosing a basis of \(\mathrm{F}'\) over F allows one to identify \(\mathrm{F}^{\prime *}\) with a subgroup of \(\mathrm{G} = \mathrm{Aut}_{\mathrm{F}}(\mathrm{F}')\) . Let \(\varphi \in \mathrm{Hom}(\mathrm{F}^{\prime *},\mathrm{K}^{*})\) ; calculate \(\mathrm{Ind}_{\mathrm{F}^{\prime *}}^{\mathrm{G}}(\varphi)\) , and prove that it is equal to \(\mathrm{Ind}_{\mathrm{F}^{\prime *}}^{\mathrm{G}}(\varphi^q)\) . Denote by \(\lambda\) the restriction of \(\varphi\) to \(\mathrm{F}^*\) , and set \(\chi_{\varphi} = \chi_{\pi_{\mathrm{P}}\otimes \pi_{\lambda ,1}} - \chi_{\pi_{\lambda ,1}} - \mathrm{Ind}_{\mathrm{F}^{\prime *}}^{\mathrm{G}}(\varphi)\) . Prove that if \(\varphi \neq \varphi^q\) , then we have \(\langle \chi_{\varphi},\chi_{\varphi}\rangle = 1\) and \(\chi_{\varphi}(1) = q - 1\) , so that \(\chi_{\varphi}\) is the character of an irreducible representation \(\pi_{\varphi}\) of dimension \(q - 1\) . Prove that we thus obtain \(\frac{1}{2} q(q - 1)\) nonisomorphic irreducible representations.

\begin{enumerate}
\def\labelenumi{\alph{enumi})}
\setcounter{enumi}{4}
\tightlist
\item
  Prove that the irreducible representations \(\pi_{\lambda}\) for \(\lambda\in\mathrm{Hom}(\mathrm{F}^{*},\mathrm{K}^{*})\) , \(\pi_{\lambda,\mu}\) for a subset \(\{\lambda,\mu\}\) of \(\mathrm{Hom}(\mathrm{F}^{*},\mathrm{K}^{*})\) with two elements, and \(\pi_{\varphi}\) for \(\varphi\in\mathrm{Hom}(\mathrm{F}^{\prime*},\mathrm{K}^{*})\) with \(\varphi\neq\varphi^{q}\) are pairwise nonisomorphic and that every irreducible representation of G is isomorphic to one of them.
\end{enumerate}

\begin{enumerate}
\def\labelenumi{\arabic{enumi})}
\setcounter{enumi}{26}
\tightlist
\item
  Let F be an algebraically closed field of characteristic p \textgreater{} 0 and G be a finite group.
\end{enumerate}

\begin{enumerate}
\def\labelenumi{\alph{enumi})}
\item
  Let \(x \in G\) . Prove that there exists a unique decomposition \(x = x_{u}x_{s}\) , where the order of \(x_{u}\) is a power of p, that of \(x_{s}\) is prime to p, and \(x_{u}\) and \(x_{s}\) commute. The elements \(x_{u}\) and \(x_{s}\) are powers of x.
\item
  Let \(\pi\) be a linear representation of \(G\) over \(F\) of finite degree. Prove that we have \(\chi_{\pi}(x) = \chi_{\pi}(x_s)\) for every \(x \in G\) .
\item
  We say that an element x of G is p-regular if its order is prime to p; we denote by \(G^{reg}\) the set of p-regular elements of G and by \(\mathcal{Z}_{\mathrm{F}}(\mathrm{G}^{\mathrm{reg}})\) the space of central functions from \(G^{reg}\) to F. From the homomorphism \(\mathrm{R}_{\mathrm{F}}(\mathrm{G}) \to \mathcal{Z}_{\mathrm{F}}(\mathrm{G}^{\mathrm{reg}})\) that sends the class of a representation \(\pi\) to the restriction of \(\chi_{\pi}\) to \(G^{reg}\) , we deduce a ring homomorphism \(\psi : F \otimes_{\mathbf{Z}} R_{\mathrm{F}}(G) \to \mathcal{Z}_{\mathrm{F}}(\mathrm{G}^{\mathrm{reg}})\) . Prove that \(\psi\) is injective (use b) and the corollary of VIII, p.~384).
\item
  Let x be a p-regular element of G, \(Z(x)\) its commutant, P a Sylow p-subgroup of \(Z(x)\) , and \(H_{x}\) the subgroup of G generated by P and x. Prove that \(H_{x}\) is the direct product of P and the subgroup of G generated by x and that its conjugacy class is well determined by the conjugacy class of x.
\item
  Construct representations \(\sigma_{1},\ldots,\sigma_{m}\) of \(H_{x}\) of degree 1 and elements \(a_{1},\ldots,a_{m}\) of F such that the function \(\sum a_{i}\chi_{\sigma_{i}}\) takes value 1 at x and 0 at the other powers of x (first construct representations of the subgroup generated by x). Let \(\rho_{i}=\operatorname{Ind}_{H_{x}}^{G}\pi_{i}\) for \(i=1,\ldots,m\) , and let \(f=\sum_{i}a_{i}\chi_{\rho_{i}}\) . Prove that \(f(x)\) is equal to \((\mathrm{Z}(x):\mathrm{H}_{x})1_{\mathrm{F}}\) , which is nonzero in F, and that \(f(y)\) is zero for every p-regular element of G that is not conjugate to x.
\end{enumerate}

\(f\) ) Deduce from \(e\) ) that the homomorphism \(\psi : \mathrm{F} \otimes_{\mathbf{Z}} \mathrm{R}_{\mathrm{F}}(\mathrm{G}) \to \mathcal{Z}_{\mathrm{F}}(\mathrm{G}^{\mathrm{reg}})\) is bijective. In particular, the number of isomorphism classes of simple F{[}G{]}-modules is equal to the number of conjugacy classes of \(p\) -regular elements of G.

\begin{enumerate}
\def\labelenumi{\arabic{enumi})}
\setcounter{enumi}{27}
\tightlist
\item
  Keep the assumptions of Exercise 27. Let \(m\) be the l.c.m. of the orders of the \(p\) -regular elements of \(G\) . Let \(F_0\) be the subfield of \(F\) generated by the \(m\) -th roots of unity; it is a finite field.
\end{enumerate}

\begin{enumerate}
\def\labelenumi{\alph{enumi})}
\item
  Prove that the character of any representation of G over F of finite degree takes all its values in \(F_{0}\) .
\item
  Prove that the quotient ring of \(F_{0}[G]\) by its radical is a product of matrix algebras over \(F_{0}\) (use a) and Wedderburn's theorem). Deduce that the canonical homomorphism \(\mathrm{R}_{\mathrm{F}_{0}}(\mathrm{G}) \to \mathrm{R}_{\mathrm{F}}(\mathrm{G})\) is bijective.
\end{enumerate}

\begin{enumerate}
\def\labelenumi{\arabic{enumi})}
\setcounter{enumi}{28}
\tightlist
\item
  Keep the assumptions and notation of Exercises 27 and 28. Choose a primitive \(m\) -th root of unity \(\zeta\) . Let \(\mathcal{O}_m\) be the ring \(\mathbf{Z}[\mathrm{X}] / (\Phi_m)\) (VIII, p.~415); we denote the class of X in \(\mathcal{O}_m\) by \(x\) . Let \(\varphi : \mathcal{O}_m \to \mathrm{F}\) be the homomorphism that sends \(x\) to \(\zeta\) . It induces an isomorphism from the subgroup of \(\mathcal{O}_m\) generated by \(x\) to the group \(\mu_m(\mathrm{F})\) ; we denote the inverse isomorphism by \(\tau\) .
\end{enumerate}

\begin{enumerate}
\def\labelenumi{\alph{enumi})}
\item
  Let \(\pi\) be a linear representation of G over F, of finite dimension \(d\) . For \(g \in \mathrm{G}\) , let \(\mathrm{P}_g(\mathrm{X}) = \prod_{i=1}^d (\mathrm{X} - \lambda_i)\) be the characteristic polynomial of \(\pi(g)\) ; denote by \(\beta_{\pi}(g)\) the element \(\sum_{i} \tau(\lambda_i)\) of \(\mathcal{O}_m\) . The resulting function \(\beta_{\pi}: \mathrm{G} \to \mathcal{O}_m\) is called the Brauer character of \(\pi\) ; we have \(\varphi \circ \beta_{\pi} = \chi_{\pi}\) and \(\beta_{\pi}(x) = \beta_{\pi}(x_s)\) for every \(x\) , so that \(\beta_{\pi}\) is determined by its restriction to \(\mathrm{G}^{\mathrm{reg}}\) .
\item
  Let \(\widehat{\mathbf{G}}\) be the set of isomorphism classes of simple F{[}G{]}-modules. Prove that the functions \((\beta_{\lambda})_{\lambda \in \widehat{\mathbf{G}}}\) are linearly independent over \(\mathbf{Z}\) . (If there exists a nontrivial relation \(\sum n_{\lambda} \beta_{\lambda} = 0\) , then we may assume that one of the \(n_{\lambda}\) is prime to \(p\) and apply \(\varphi\) to arrive at a contradiction using Exercise 27, f.)
\item
  Let \(\mathcal{B}\) be the set of functions from \(\mathrm{G}^{\mathrm{reg}}\) to \(\mathcal{O}_m\) ; the homomorphism \(\mathrm{R}_{\mathrm{F}}(\mathrm{G}) \to \mathcal{B}\) that sends the class of \(\pi\) to \(\beta_{\pi}\) is injective. Deduce that two semisimple F{[}G{]}-modules are isomorphic if and only if their Brauer characters are equal.
\end{enumerate}

\begin{enumerate}
\def\labelenumi{\arabic{enumi})}
\setcounter{enumi}{29}
\tightlist
\item
  Let \(\mathrm{F}\) be an algebraically closed field of characteristic \(p > 0\) , let \(\mathbf{F}_p\) be the prime subfield of \(\mathrm{F}\) , and let \(\mathrm{G} = \mathbf{SL}_2(\mathbf{F}_p)\) . Denote by \(\mathrm{V}\) the F-vector space \(\mathrm{F} \otimes_{\mathbf{F}_p} \mathbf{F}_p^2\) endowed with the action of \(\mathrm{G}\) deduced from the action on \(\mathbf{F}_p^2\) .
\end{enumerate}

\begin{enumerate}
\def\labelenumi{\alph{enumi})}
\item
  Prove that the F{[}G{]}-module \(\mathbf{S}^{d}(V)\) is simple for \(0 \leqslant d < p\) . (Let W be a nonzero F{[}G{]}-submodule of \(\mathbf{S}^{d}(V)\) , and let \((e, f)\) be a basis of V. Prove that \(e^{d}\) belongs to W using the matrix \(\begin{pmatrix}1 & 1 \\ 0 & 1\end{pmatrix}\) , and then that W is equal to \(\mathbf{S}^{d}(V)\) using the matrix \(\begin{pmatrix}1 & 0 \\ 1 & 1\end{pmatrix}\) .) b) Prove that the F{[}G{]}-module \(\mathbf{S}^{p}(V)\) is not simple (consider the image of V by the mapping \(v \mapsto v^{p}\) ).
\setcounter{enumi}{2}
\item
  For \(p \geqslant 7\) , prove that the dimension of \(\mathbf{S}^{p-3}(\mathrm{V})\) does not divide the order of G. d) Prove that every simple F{[}G{]}-module is isomorphic to one of the modules \(\mathbf{S}^{d}(\mathrm{V})\) for \(0 \leqslant d < p\) (use Exercise 27, f)).
\end{enumerate}

\appendix
\renewcommand{\thechapter}{\arabic{chapter}}
\chapter{Algebras without Unit Element}

In this appendix, k is a commutative ring; the k-algebras are assumed associative but not necessarily unital.

\section{Regular Ideals}

Let A be a k-algebra. Recall (III, §1, No.~2, p.~430) that a left ideal of A is a k-submodule \(\mathfrak{a}\) of A such that the relations \(a \in A\) , \(x \in \mathfrak{a}\) imply \(ax \in \mathfrak{a}\) . We define the notions of right ideal and two-sided ideal likewise. If A is a k-algebra and \(\mathfrak{a}\) a two-sided ideal of A, then when passing to the quotients, the multiplication in A defines a k-algebra structure on the k-module A/\(\mathfrak{a}\) (loc. cit.).

A left ideal of A is called maximal if it is a maximal element of the set of proper left ideals of A for the inclusion.

DEFINITION 1. --- Let A be a k-algebra and \(\mathfrak{a}\) a left ideal of A. An element u of A such that au - a belongs to \(\mathfrak{a}\) for every \(a \in A\) is called a right unit modulo \(\mathfrak{a}\) . We say that the ideal \(\mathfrak{a}\) is regular if there exists a right unit modulo \(\mathfrak{a}\) .

Likewise, we say that a right ideal \(\mathfrak{b}\) of A is regular if A has a left unit modulo \(\mathfrak{b}\) , that is, an element v such that \(va - a \in \mathfrak{b}\) for every \(a \in A\) . When an ideal is two-sided, one needs to specify whether it is regular as a left ideal or as a right ideal.

When the algebra A is unital, the unit element of A is a right unit modulo \(\mathfrak{a}\) for every left ideal \(\mathfrak{a}\) of A; in this case, every left (or right) ideal of A is regular. On the other hand, if A is a k-algebra whose elements are all nilpotent, then A is the only (left or right) ideal of A that is regular.

Let \(\mathfrak{a}\) be a regular left ideal of A and \(u\) a right unit modulo \(\mathfrak{a}\) . If \(\mathfrak{b}\) is a left ideal of A containing \(\mathfrak{a}\) , then \(u\) is a right unit modulo \(\mathfrak{a}\) , so \(\mathfrak{b}\) is regular. So the left ideals of A that are maximal and regular are the maximal elements of the set of regular proper left ideals of A. Moreover, a left ideal \(\mathfrak{b}\) of A containing \(\mathfrak{a}\) is proper if and only if \(u\) does not belong to \(\mathfrak{b}\) . The set of proper left ideals \(\mathfrak{b}\) of A containing \(\mathfrak{a}\) , ordered by inclusion, is therefore an inductive set. Theorem 2 of Set Theory, III, §2, No.~4, p.~154 applied to this set gives the following result.

PROPOSITION 1. --- Every regular proper left (resp. right) ideal of A is contained in a regular maximal left (resp. right) ideal.

PROPOSITION 2. --- Let A be a commutative k-algebra. An ideal \(\mathfrak{a}\) of A is a regular maximal ideal if and only if the pseudoring A/ \(\mathfrak{a}\) (I, §8, No.~1, p.~98) is a field.

An element u of A is a (right or left) unit modulo \(\mathfrak{a}\) if and only if the canonical image of u in A/\(\mathfrak{a}\) is a unit element of the algebra A/\(\mathfrak{a}\) . Therefore, the ideal \(\mathfrak{a}\) is regular if and only if the algebra A/\(\mathfrak{a}\) is unital. Suppose that these conditions are satisfied.

The mapping \(\mathfrak{b} \mapsto \mathfrak{b}/\mathfrak{a}\) is a bijection from the set of ideals of the algebra A containing \(\mathfrak{a}\) to the set of ideals of the algebra A/\(\mathfrak{a}\) . These are the ideals of the ring A/\(\mathfrak{a}\) because this algebra is unital. Finally, by Theorem 1 of I, §9, No.~1, p.~115, the ring A/\(\mathfrak{a}\) is a field if and only if it is nonzero and its only ideals are 0 and itself. The proposition follows.

Examples. --- 1) Let V be an infinite-dimensional vector space over a commutative field K. Let \(\operatorname{End}_{\mathrm{K}}^{\mathrm{f}}(\mathrm{V})\) be the K-subalgebra of \(\operatorname{End}_{\mathrm{K}}(\mathrm{V})\) consisting of the endomorphisms of finite rank. Let W be a linear subspace of V, and let \(\mathfrak{a}_{W}\) be the set of elements of \(\operatorname{End}_{\mathrm{K}}^{\mathrm{f}}(\mathrm{V})\) whose kernel contains W; it is a left ideal of \(\operatorname{End}_{\mathrm{K}}^{\mathrm{f}}(\mathrm{V})\) . An element u of \(\operatorname{End}_{\mathrm{K}}^{\mathrm{f}}(\mathrm{V})\) is a right unit modulo \(\mathfrak{a}_{W}\) if and only if we have \(u(x)=x\) for every \(x\in W\) . Such an element u exists, that is, \(\mathfrak{a}_{W}\) is regular, if and only if W is finite-dimensional. The ideal \(\mathfrak{a}_{W}\) is maximal and regular if and only if W has dimension 1.

*2) Let \(\mathrm{T}\) be a locally compact space, and let \(\mathcal{C}_0(\mathrm{T})\) be the commutative \(\mathbf{C}\) -algebra of continuous mappings from \(\mathrm{T}\) to \(\mathbf{C}\) that tend to 0 at infinity (Gen.~Top., X, §4, No.~4, p.~316); it is unital if and only if \(\mathrm{T}\) is compact. Let \(\mathrm{F}\) be a closed subset of \(\mathrm{T}\) , and let \(\mathfrak{a}_{\mathrm{F}}\) be the set of elements of \(\mathcal{C}_0(\mathrm{T})\) whose restriction to F is the zero function; it is an ideal of \(\mathcal{C}_{0}(T)\) . An element u of \(\mathcal{C}_{0}(T)\) is a unit modulo \(\mathfrak{a}_{F}\) if and only if we have \(u(t)=1\) for every \(t\in F\) . Such an element u exists, that is, \(\mathfrak{a}_{F}\) is regular, if and only if F is compact. The mapping \(t\mapsto \mathfrak{a}_{\{t\}}\) is a bijection from T to the set of regular maximal ideals of \(\mathcal{C}_{0}(T)\) (TS, I, §3, n° 2, p.~32, corollaire 1). Suppose that T is not compact, and denote by \(\mathfrak{a}\) the subset of \(\mathcal{C}_{0}(T)\) consisting of the functions with compact support; then \(\mathfrak{a}\) is an ideal of \(\mathcal{C}_{0}(T)\) that is not contained in any regular ideal of \(\mathcal{C}_{0}(T)\) .

\begin{enumerate}
\def\labelenumi{\arabic{enumi})}
\setcounter{enumi}{2}
\tightlist
\item
  Let \(L^{1}(\mathbf{R})\) be the convolution algebra of the locally compact group R. Recall (cf.~Int., VIII, §4, No.~5, p.~38) that \(L^{1}(\mathbf{R})\) is the space of classes of functions on R that are integrable for the Lebesgue measure. The product of the classes of two functions f and g is the class of the function \(f * g\) defined by the formula
\end{enumerate}

\[
(f * g) (s) = \int_ {- \infty} ^ {+ \infty} f (t) g (s - t) d t
\]

for almost every \(s \in \mathbf{R}\) . The algebra \(\mathrm{L}^1(\mathbf{R})\) is not unital. For any \(a\) in \(\mathbf{R}\) , denote by \(\mathfrak{m}_a\) the set of elements \(f\) of \(\mathrm{L}^1(\mathbf{R})\) satisfying

\[
\int_ {- \infty} ^ {+ \infty} f (t) e ^ {- i a t} d t = 0.
\]

By TS, II, §3, n° 2, p.~252, théorème 1, the mapping \(a \mapsto \mathfrak{m}_a\) is a bijection from \(\mathbf{R}\) to the set of regular maximal ideals of the algebra \(\mathrm{L}^1(\mathbf{R})\) .*

\section{Adjunction of a Unit Element}

Let A be a k-algebra. In III, §1, No.~2, p.~431, we defined the unital algebra \(\widetilde{A}\) deduced from A by the adjunction of a unit element e. We identify A with a two-sided ideal of \(\widetilde{A}\) . The k-module \(\widetilde{A}\) is the direct sum of the submodules ke and A.

PROPOSITION 3. --- a) Let \(\widetilde{\mathfrak{a}}\) be a left ideal of \(\widetilde{\mathrm{A}}\) such that \(\widetilde{\mathrm{A}} = \widetilde{\mathfrak{a}} + \mathrm{A}\) . We set \(\mathfrak{a} = \widetilde{\mathfrak{a}} \cap \mathrm{A}\) . There exists an element \(u\) of \(\mathrm{A}\) such that \(u - e\) belongs to \(\widetilde{\mathfrak{a}}\) . If \(u\) is such an element, then it is a right unit of \(\mathrm{A}\) modulo \(\mathfrak{a}\) , and we have \(\widetilde{\mathfrak{a}} = \mathfrak{a} + k(u - e)\) ; in particular, the ideal \(\mathfrak{a}\) is regular.

\begin{enumerate}
\def\labelenumi{\alph{enumi})}
\setcounter{enumi}{1}
\item
  Conversely, let \(\mathfrak{a}\) be a regular left ideal of A and \(u\) a right unit of A modulo \(\mathfrak{a}\) . Set \(\widetilde{\mathfrak{a}} = \mathfrak{a} + k(u - e)\) . Then \(\widetilde{\mathfrak{a}}\) is a left ideal of \(\widetilde{\mathrm{A}}\) such that \(\widetilde{\mathrm{A}} = \widetilde{\mathfrak{a}} + \mathrm{A}\) , and we have \(\mathfrak{a} = \widetilde{\mathfrak{a}} \cap \mathrm{A}\) .
\item
  If \(k\) is a field, then the condition \(\widetilde{\mathrm{A}} = \widetilde{\mathfrak{a}} + \mathrm{A}\) is equivalent to saying that \(\widetilde{\mathfrak{a}}\) is not contained in \(\mathrm{A}\) .
\end{enumerate}

Under the assumption of a), the element \(e\) of \(\widetilde{\mathbf{A}}\) can be written as \((e - u) + u\) with \(u \in \mathbf{A}\) and \(u - e \in \widetilde{\mathfrak{a}}\) . Let \(x = \lambda e + a\) be an element of \(\widetilde{\mathbf{A}}\) , where \(\lambda \in k\) and \(a \in \mathbf{A}\) ; if \(x\) belongs to \(\widetilde{\mathfrak{a}}\) , then the element \(y = a + \lambda u = x + \lambda (u - e)\) of \(\mathbf{A}\) belongs to \(\mathfrak{a}\) , and we have \(x = y - \lambda (u - e)\) ; this proves the equality \(\widetilde{\mathfrak{a}} = \mathfrak{a} + k(u - e)\) . Moreover, for every \(a\) in \(\mathbf{A}\) , the element \(au - a = a(u - e)\) of \(\mathbf{A}\) belongs to \(\mathfrak{a}\) , so \(u\) is a right unit modulo \(\mathfrak{a}\) . This proves assertion a).

Let \(\mathfrak{a}\) and \(u\) be as in b). Since \(\widetilde{\mathrm{A}}\) is the direct sum of \(\mathrm{A}\) and \(k(u - e)\) , we have \(\widetilde{\mathfrak{a}} + \mathrm{A} = \widetilde{\mathrm{A}}\) and \(\widetilde{\mathfrak{a}} \cap \mathrm{A} = \mathfrak{a}\) . Every element of \(\widetilde{\mathfrak{a}}\) is of the form \(x + \lambda (u - e)\) with \(x \in \mathfrak{a}\) and \(\lambda \in k\) . For every \(a \in \mathrm{A}\) , we have

\[
a (x + \lambda (u - e)) = a x + \lambda (a u - a).
\]

Since u is a right unit of A modulo \(\mathfrak{a}\), the sum \(ax + \lambda(au - a)\) belongs to \(\mathfrak{a}\), so that \(\widetilde{\mathfrak{a}}\) is a left ideal of A. This proves b).

Finally, if \(k\) is a field, then \(A\) is a hyperplane in \(\widetilde{A}\) , and \(\widetilde{A}\) is the sum of \(\widetilde{\mathfrak{a}}\) and \(A\) if and only if \(\widetilde{\mathfrak{a}}\) is not contained in \(A\) .

COROLLARY. --- The regular left ideals of A are the ideals of the form \(A \cap \widetilde{\mathfrak{a}}\) , where \(\widetilde{\mathfrak{a}}\) is a left ideal of \(\widetilde{A}\) such that \(\widetilde{A} = \widetilde{\mathfrak{a}} + A\) .

PROPOSITION 4. --- a) Let \(\mathfrak{a}\) be a regular maximal left ideal of A. There exists a unique left ideal \(\widetilde{\mathfrak{a}}\) of \(\widetilde{\mathrm{A}}\) such that \(\widetilde{\mathrm{A}} = \widetilde{\mathfrak{a}} + \mathrm{A}\) and \(\mathfrak{a} = \widetilde{\mathfrak{a}} \cap \mathrm{A}\) . This ideal is maximal and does not contain A.

\begin{enumerate}
\def\labelenumi{\alph{enumi})}
\setcounter{enumi}{1}
\tightlist
\item
  The mapping \(\widetilde{\mathfrak{a}}\mapsto \widetilde{\mathfrak{a}}\cap \mathrm{A}\) is a bijection from the set of maximal left ideals of \(\widetilde{\mathrm{A}}\) not containing A to the set of regular maximal left ideals of A.
\end{enumerate}

Let \(\mathfrak{a}\) be a regular maximal left ideal of A. By Proposition 3, a), the left ideals \(\mathfrak{b}\) of \(\widetilde{\mathrm{A}}\) such that \(\mathfrak{b} + \mathrm{A} = \widetilde{\mathrm{A}}\) and \(\mathfrak{b} \cap \mathrm{A} = \mathfrak{a}\) are the ideals \(\mathfrak{a} + k(u - e)\) , where \(u\) is a right unit modulo \(\mathfrak{a}\) . To prove the uniqueness of \(\widetilde{\mathfrak{a}}\) , it therefore suffices to prove that two right units \(u\) and \(u'\) of A modulo \(\mathfrak{a}\) are congruent modulo \(\mathfrak{a}\) . Let us reason by contradiction and suppose that \(u - u'\) does not belong to \(\mathfrak{a}\) . The formula \(x(u - u') = (xu - x) - (xu' - x)\) shows that we have \(\mathrm{A}(u - u') \subset \mathfrak{a}\) ; it follows that \(\mathfrak{a} + k(u - u')\) is a left ideal of A containing \(\mathfrak{a}\) and distinct from \(\mathfrak{a}\) . Since \(\mathfrak{a}\) is maximal, we therefore have \(\mathfrak{a} + k(u - u') = \mathrm{A}\) , and so \(\mathrm{AA} \subset \mathfrak{a}\) . For every \(x \in \mathrm{A}\) , we have \(x \equiv xu (\mathrm{mod} \mathfrak{a})\) , and therefore \(x \in \mathfrak{a}\) by the above, which contradicts the assumption \(\mathfrak{a} \neq \mathrm{A}\) .

Hence there exists a unique left ideal \(\widetilde{\mathfrak{a}}\) of \(\widetilde{\mathrm{A}}\) such that \(\widetilde{\mathrm{A}} = \widetilde{\mathfrak{a}} + \mathrm{A}\) and \(\mathfrak{a} = \widetilde{\mathfrak{a}} \cap \mathrm{A}\) . Let \(\mathfrak{b}\) be a proper left ideal of \(\widetilde{\mathrm{A}}\) containing \(\widetilde{\mathfrak{a}}\) . Then \(\mathfrak{b} \cap \mathrm{A}\) is a proper left ideal of \(\mathrm{A}\) containing \(\mathfrak{a}\) . It is therefore equal to \(\mathfrak{a}\) because \(\mathfrak{a}\) is maximal, which implies \(\mathfrak{b} = \widetilde{\mathfrak{a}}\) by Lemma 2 of VIII, p.~4. This proves that \(\widetilde{\mathfrak{a}}\) is a maximal ideal of \(\widetilde{\mathrm{A}}\) ; for such an ideal, the condition \(\widetilde{\mathrm{A}} = \widetilde{\mathfrak{a}} + \mathrm{A}\) means that \(\widetilde{\mathfrak{a}}\) does not contain A.

It remains to prove that if \(\widetilde{\mathfrak{a}}\) is a maximal left ideal of \(\tilde{A}\) that does not contain A, then the left ideal \(\mathfrak{a} = A \cap \widetilde{\mathfrak{a}}\) of A is maximal. Let \(\mathfrak{b}\) be a proper left ideal of A containing \(\mathfrak{a}\). Let u be a right unit modulo \(\mathfrak{a}\) such that \(\widetilde{\mathfrak{a}} = \mathfrak{a} + k(u - e)\) (Proposition 3). It is also a right unit modulo \(\mathfrak{b}\); denote by \(\widetilde{\mathfrak{b}}\) the ideal \(\mathfrak{b} + k(u - e)\) of \(\tilde{A}\) . By Proposition 3, b), we have \(\mathfrak{b} = A \cap \widetilde{\mathfrak{b}}\) . The ideal \(\widetilde{\mathfrak{a}}\) , equal to \(\mathfrak{a} + k(u - e)\) , is contained in \(\widetilde{\mathfrak{b}}\) and is therefore equal to \(\widetilde{\mathfrak{b}}\) because \(\widetilde{\mathfrak{b}}\) is distinct from \(\tilde{A}\) and \(\widetilde{\mathfrak{a}}\) is maximal. Consequently, we have \(\mathfrak{a} = \mathfrak{b}\) , and \(\mathfrak{a}\) is maximal.

We leave it to the reader to translate Propositions 3 and 4 to right ideals.

\section{The Radical of an Algebra}

Let A be a k-algebra. A left pseudomodule over A is a k-module M endowed with the structure of a left pseudomodule over the pseudoring A (II, Appendix, No.~2, p.~378) such that we have \(a(\lambda x) = \lambda(ax) = (\lambda a)x\) for \(\lambda \in k\) , \(a \in A\) , and \(x \in M\) . We define right pseudomodules over A likewise.

Let M be a left pseudomodule over A; we define a left \(\tilde{A}\) -module structure on M called canonical by setting

\[
(\lambda e + a) x = \lambda x + a x \quad \text { for } \lambda \in k, a \in A, x \in M.
\]

Conversely, every left \(\widetilde{A}\) -module is canonically endowed, by restricting the ring of operators to k and A, with the structure of a left pseudomodule over A. Thus, the species of structures of left pseudomodules over A and of modules over \(\widetilde{A}\) are equivalent (Set Theory, IV, §1, No.~7, p.~262).

Let M be a left pseudomodule over A, and let N be a k-submodule of M. Then N is an \(\widetilde{A}\) -submodule of M if and only if it is stable under the action of A; we then say that N is a sub-pseudomodule of M.

As in the case of rings, we define the left pseudomodule \(A_{s}\) over A and the right pseudomodule \(A_{d}\) . The left (resp. right) ideals of A are the sub-pseudomodules of \(A_{s}\) (resp. \(A_{d}\) ).

Let \(\mathfrak{a}\) be a regular left ideal of A and \(u\) a right unit modulo \(\mathfrak{a}\) . Set \(\mathrm{M} = \mathrm{A}_s / \mathfrak{a}\) , and denote the image of \(u\) in M by \(z\) . We have \(\mathrm{M} = \mathrm{A}z\) , and \(\mathfrak{a}\) is the annihilator of \(z\) .

Conversely, let M be a left pseudomodule over A, and let z be an element of M such that M = Az. Then there exists an element u of A such that z = uz. For every \(a \in A\) , we have \((au - a)z = 0\) , so au - a belongs to the annihilator \(\mathfrak{a}\) of z. Consequently, \(\mathfrak{a}\) is a regular left ideal of A, the element u is a right unit modulo \(\mathfrak{a}\), and when passing to the quotient, the mapping \(a \mapsto az\) defines an isomorphism from \(A_{s}/\mathfrak{a}\) to M.

DEFINITION 2. --- We say that a pseudomodule M over A is simple if we have AM ≠ 0 and if 0 and M are the only sub-pseudomodules of M.

This corresponds to saying that the \(\tilde{A}\) -module M is simple and that its annihilator does not contain A. When A is a ring and M is an A-module, we have AM = M, so Definition 2 coincides with Definition 1 of VIII, p.~45.

PROPOSITION 5. --- a) Let \(\mathfrak{m}\) be a regular maximal left ideal of A. Then the pseudomodule \(A_s / \mathfrak{m}\) is simple, and there exists a nonzero element of \(A_s / \mathfrak{m}\) with annihilator \(\mathfrak{m}\) .

\begin{enumerate}
\def\labelenumi{\alph{enumi})}
\setcounter{enumi}{1}
\item
  Let M be a simple pseudomodule, and let x be a nonzero element of M. We have M = Ax, the annihilator \(\mathfrak{m}\) of x is a regular maximal left ideal of A, and when passing to the quotient, the mapping \(a \mapsto ax\) defines an isomorphism from \(A_{s}/\mathfrak{m}\) to M.
\item
  Let M be a pseudomodule not reduced to 0. Then M is simple if and only if we have M = Ax for every nonzero element x of M.
\end{enumerate}

Let \(\mathfrak{m}\) be a regular maximal left ideal of A. We have seen that there exists a nonzero element \(z\) of \(\mathrm{A}_s / \mathfrak{m}\) with annihilator equal to \(\mathfrak{m}\) and such that \(\mathrm{A}z = \mathrm{A}_s / \mathfrak{m}\) ; in particular, we have \(\mathrm{A}(\mathrm{A}_s / \mathfrak{m}) \neq 0\) . Moreover, every sub-pseudomodule of \(\mathrm{A}_s / \mathfrak{m}\) is of the form \(\mathfrak{n} / \mathfrak{m}\) , where \(\mathfrak{n}\) is a left ideal of A containing \(\mathfrak{m}\) . Since \(\mathfrak{m}\) is maximal, the only possibilities are \(\mathfrak{n} = \mathfrak{m}\) and \(\mathfrak{n} = \mathrm{A}\) , so that the pseudomodule \(\mathrm{A}_s / \mathfrak{a}\) is simple. This proves a).

Under the assumptions of b), the set of elements y of M such that Ay = 0 is a sub-pseudomodule of M different from M and therefore reduced to 0. Consequently, Ax is a nonzero sub-pseudomodule of M, which implies M = Ax. Assertion b) then follows from the remarks before Definition 2.

Let M be a nonzero pseudomodule. Suppose that we have Ax = M for every nonzero element x of M. In particular, we have AM ≠ 0. Let N be a nonzero sub-pseudomodule of M, and let x be a nonzero element of N. We have Ax = M, and therefore N = M. Hence M is simple. Conversely, if M is simple, then we have M = Ax for every \(x \neq 0\) by b).

DEFINITION 3. --- The radical of the k-algebra A, denoted by \(\Re(\mathrm{A})\) , is the intersection of the regular maximal left ideals of A.

When A is a ring, every left ideal of A is regular, so the definition of the radical coincides with Definition 2 of VIII, p.~154.

Examples. --- 1) The radical of the \(k\) -algebra \(\operatorname{End}_k^{\mathrm{f}}(\mathrm{V})\) is reduced to 0 (VIII, p.~436, Example 1). *The same holds for the algebras \(\mathcal{C}_0(\mathrm{T})\) and \(\mathrm{L}^1 (\mathbf{R})\)).*

\begin{enumerate}
\def\labelenumi{\arabic{enumi})}
\setcounter{enumi}{1}
\tightlist
\item
  Let A be a pseudoring in which all elements are nilpotent. The radical of A is equal to A because A is the only regular left ideal.
\end{enumerate}

Proposition 5 immediately implies the following result.

PROPOSITION 6. --- The radical of the algebra A is the intersection of the annihilators of the simple pseudomodules. It is, in particular, a two-sided ideal of A.

PROPOSITION 7. --- The radical of A is the trace on A of the radical of \(\widetilde{A}\) ; it is also equal to the radical of the opposite algebra \(A^{\circ}\) (that is, to the intersection of the regular maximal right ideals of A). If the ring k is without radical, then the radical of A is equal to that of \(\widetilde{A}\) .

The equality \(\Re (\widetilde{\mathrm{A}})\cap \mathrm{A} = \Re (\mathrm{A})\) follows from Proposition 4 of VIII, p.~438, b). Since \(\widetilde{\mathrm{A}}\) and \(\widetilde{\mathrm{A}}^{\circ}\) have the same radical (VIII, p.~156, Corollary 1), the equality \(\Re (\mathrm{A}) = \Re (\mathrm{A}^{\circ})\) follows. If \(k\) is without radical, then the intersection of the maximal left ideals of \(\widetilde{\mathrm{A}}\) containing A is equal to A. Consequently, \(\Re (\widetilde{\mathrm{A}})\) is contained in A and therefore equal to \(\Re (\mathrm{A})\) .

Remark. --- Let x and y be elements of A; we say that x is a left adverse of y, or that y is a right adverse of x, if, in \(\widetilde{A}\) , x - e is a left inverse of y - e, that is, if we have \(x + y = xy\) . By Proposition 7 and Jacobson's theorem (VIII, p.~156, Theorem 1), the radical of A consists of the elements x of A such that ux - e is left invertible in \(\widetilde{A}\) for every u in \(\widetilde{A}\) . Since we have \((a + \lambda e)x - e = ax + \lambda x - e\) (for \(a \in A\) , \(\lambda \in k\) ), the radical of A is the set of elements x of A such that \(ax + \lambda x\) has a left adverse in A for every \(a \in A\) and \(\lambda \in k\) .

\section{Density Theorem}

Let A be a k-algebra. A left pseudomodule M is called semisimple if it is the direct sum of a family of simple left pseudomodules.

Lemma 1. --- Let M be a semisimple left A-pseudomodule. Let B be the bicommutant of the \(\widetilde{A}\) -module M. Then every A-sub-pseudomodule of M is a B-submodule of M.

The \(\widetilde{\mathrm{A}}\) -module M is semisimple. Let N be an A-sub-pseudomodule of M. Then N is an \(\widetilde{\mathrm{A}}\) -submodule of M. By Corollary 2 of VIII, p.~56, there exists a projector \(p\) of the A-module M with image N. Since we have the relation \(pb = bp\) for every \(b \in \mathrm{B}\) , we obtain that N is a B-submodule of M.

Lemma 2. --- Let M be a semisimple left A-pseudomodule, and let x be an element of M. There exists an element \(a \in A\) such that ax = x.

Let N be the left A-pseudomodule \(\widetilde{Ax}/Ax\) . It satisfies \(AN = \{0\}\) . By Corollary 3 of VIII, p.~56 applied to the \(\widetilde{A}\) -modules M and \(\widetilde{Ax}\) , the A-pseudomodule N is semisimple. By definition, every simple sub-pseudomodule S of N satisfies \(AS \neq \{0\}\) . Consequently, the pseudomodule N is zero, and we obtain that \(x \in Ax\) .

THEOREM 1 (Jacobson's density theorem). --- Let M be a semisimple left A-pseudomodule. Let b be an element of the bicommutant of the \(\widetilde{\mathrm{A}}\) -module M. Let \(\mathrm{F} = \{x_1, \ldots, x_n\}\) be a finite subset of M. Then there exists an element \(a \in \mathrm{A}\) such that \(bx_i = ax_i\) for every \(i \in [1, n]\) .

Let B be the bicommutant of the \(\widetilde{\mathrm{A}}\) -module M. The A-pseudomodule \(\mathrm{M}^n\) is semisimple. Let \(\boldsymbol{x} = (x_1, \ldots, x_n) \in \mathrm{M}^n\) . It follows from Lemma 2 that \(\boldsymbol{x} \in \mathrm{A}\boldsymbol{x}\) . The bicommutant of the \(\widetilde{\mathrm{A}}\) -module \(\mathrm{M}^n\) coincides with the homotheties of the B-module \(\mathrm{M}^n\) (VIII, p.~79, Proposition 2). By Lemma 1, the A-subpseudomodule \(\mathrm{A}\boldsymbol{x}\) of \(\mathrm{M}^n\) is therefore a B-submodule of \(\mathrm{M}^n\) . We therefore have the inclusion \(\mathrm{B}\boldsymbol{x} \subset \mathrm{A}\boldsymbol{x}\) . So there exists an \(a \in \mathrm{A}\) such that \(b\boldsymbol{x} = a\boldsymbol{x}\) . The result follows.
